\documentclass[11pt,a4paper]{article}
% Shared preamble. Kept separate so the figure-producing sibling can reuse it.
\usepackage[T1]{fontenc}
\usepackage[utf8]{inputenc}
\usepackage{lmodern}
\usepackage[a4paper,margin=2.6cm]{geometry}
\usepackage{amsmath}
\usepackage{amssymb}
\usepackage{booktabs}
\usepackage{longtable}
\usepackage{array}
\usepackage{graphicx}
\usepackage{xcolor}
\usepackage[font=small,labelfont=bf]{caption}
\usepackage{enumitem}
\usepackage{fancyhdr}
\usepackage{microtype}
\usepackage[hidelinks,breaklinks]{hyperref}

\definecolor{boxrule}{HTML}{9AA3AF}
\definecolor{boxfill}{HTML}{F4F5F7}
\definecolor{warnfill}{HTML}{FDF3F2}
\definecolor{warnrule}{HTML}{C0563F}

\setlength{\parskip}{0.45em}
\setlength{\parindent}{0pt}
\renewcommand{\arraystretch}{1.18}

\pagestyle{fancy}
\fancyhf{}
\fancyhead[L]{\small Benchmarking the GT screener}
\fancyhead[R]{\small\thepage}
\renewcommand{\headrulewidth}{0.3pt}

% ---------------------------------------------------------------------------
% Callout boxes
% ---------------------------------------------------------------------------
\newcommand{\boundary}[1]{%
  \par\medskip\noindent
  \fcolorbox{warnrule}{warnfill}{%
    \begin{minipage}{0.965\linewidth}\small\vspace{2pt}
    \textbf{Claim boundary.}\ #1\vspace{2pt}
    \end{minipage}}\par\medskip}

\newcommand{\assumption}[1]{%
  \par\medskip\noindent
  \fcolorbox{boxrule}{boxfill}{%
    \begin{minipage}{0.965\linewidth}\small\vspace{2pt}
    \textbf{Assumption to confirm.}\ #1\vspace{2pt}
    \end{minipage}}\par\medskip}

\newcommand{\simnote}[1]{%
  \par\medskip\noindent
  \fcolorbox{warnrule}{warnfill}{%
    \begin{minipage}{0.965\linewidth}\small\vspace{2pt}
    \textbf{Simulated data.}\ #1\vspace{2pt}
    \end{minipage}}\par\medskip}

\newcommand{\keyfinding}[1]{%
  \par\medskip\noindent
  \fcolorbox{boxrule}{boxfill}{%
    \begin{minipage}{0.965\linewidth}\vspace{2pt}
    #1\vspace{2pt}
    \end{minipage}}\par\medskip}

% ---------------------------------------------------------------------------
% Include a generated table body inside a tabular.
%
% LaTeX's \input{} appends \relax after the file, which lands after the final
% \\ and puts the following \bottomrule's \noalign inside a cell ("Misplaced
% \noalign").  The TeX primitive does not, so use it for table bodies.
% ---------------------------------------------------------------------------
\makeatletter
\newcommand{\tabinput}[1]{\@@input #1 }
\makeatother

% ---------------------------------------------------------------------------
% Figure placeholders. The sibling analysis drops PDFs into figures/ ; until
% then a labelled frame stands in, so the document always builds.
% #1 file stem  #2 label  #3 caption
% ---------------------------------------------------------------------------
\newcommand{\figslot}[3]{%
  \begin{figure}[htbp]
  \centering
  \IfFileExists{figures/#1.pdf}
    {\includegraphics[width=0.92\linewidth]{figures/#1}}
    {\fcolorbox{boxrule}{boxfill}{%
       \begin{minipage}[c][4.6cm][c]{0.9\linewidth}
       \centering\small
       \textbf{[ figure pending ]}\par\smallskip
       \texttt{figures/#1.pdf}\par\smallskip
       Produced by the companion synthetic-data analysis.\\
       The surrounding argument does not depend on its values.
       \end{minipage}}}
  \caption{#3}
  \label{#2}
  \end{figure}}

% shorthand
\newcommand{\dR}{\ensuremath{\Delta R^{2}}}
\newcommand{\Rsq}{\ensuremath{R^{2}}}

% GENERATED BY scripts/compute.py -- DO NOT EDIT BY HAND
\newcommand{\NcampusCohort}{46}
\newcommand{\NanywhereCohort}{300}
\newcommand{\Nbothcohorts}{346}
\newcommand{\RelComparator}{0.94}
\newcommand{\RelCompLow}{0.90}
\newcommand{\RelCompHigh}{0.98}
\newcommand{\RelMAP}{0.95}
\newcommand{\RelScreener}{0.90}
\newcommand{\RhoCogATobserved}{0.63}
\newcommand{\RhoCogAToperational}{0.31}
\newcommand{\BlockAdministered}{30}
\newcommand{\UniformAttenR}{0.9772}
\newcommand{\UniformAttenRsq}{0.9549}
\newcommand{\UniformAttenLossPct}{4.5\%}
\newcommand{\MDEcampus}{0.40}
\newcommand{\MDEboth}{0.15}
\newcommand{\CIcampusLo}{0.07}
\newcommand{\CIcampusHi}{0.58}
\newcommand{\CIbothLo}{0.25}
\newcommand{\CIbothHi}{0.44}
\newcommand{\NforThirtyFive}{62}
\newcommand{\NciTenAtThirtyFive}{298}
\newcommand{\NciSevenFiveAtThirtyFive}{528}
\newcommand{\DEcampus}{3.25}
\newcommand{\EffNcampus}{14}
\newcommand{\NincrThreeMid}{191}
\newcommand{\NincrFiveMid}{112}
\newcommand{\SpuriousMid}{0.0043}
\newcommand{\SpuriousHigh}{0.0084}
\newcommand{\SpuriousLow}{0.0003}
\newcommand{\SpuriousWorst}{0.0188}
\newcommand{\NdetectSpurious}{1376}
\newcommand{\NdetectSpuriousHigh}{700}
\newcommand{\NshiftedFive}{172}
\newcommand{\NshiftedThree}{380}
\newcommand{\PPVFortyTen}{0.27}
\newcommand{\PPVFortyTenPct}{27\%}
\newcommand{\FPFortyTenPct}{73\%}
\newcommand{\PPVliftFortyTen}{2.7}
\newcommand{\PPVtwoPercentGood}{0.27}
\newcommand{\PPVtwoPercentGoodPct}{27\%}
\newcommand{\PPVfivePercentGoodPct}{49\%}
\newcommand{\PPVtwentyGoodPct}{82\%}
\newcommand{\SpecForHalfPPVatFive}{0.953}
\newcommand{\SpecForHalfPPVatTwenty}{0.775}
\newcommand{\DCninetyTen}{0.94}
\newcommand{\DAninetyTen}{0.96}
\newcommand{\DisagreeNinetyTenPct}{6\%}
\newcommand{\DCninetyFifty}{0.86}
\newcommand{\RestrictedTenFromForty}{0.18}
\newcommand{\AttenuationTenPct}{56\%}
\newcommand{\NunrestrictedForty}{47}
\newcommand{\NrestrictedForty}{250}
\newcommand{\NmultiplierTen}{5.3}
\newcommand{\SDratioTen}{0.411}
\newcommand{\NmultiplierTwenty}{4.1}
\newcommand{\LRdeltaThirty}{0.007}
\newcommand{\LRdeltaSixty}{0.015}
\newcommand{\LRseThirty}{0.063}
\newcommand{\LRseSixty}{0.028}
\newcommand{\LRfloorThirty}{0.0097}
\newcommand{\LRfloorSixty}{0.0094}
\newcommand{\LRseDropPct}{56\%}
\newcommand{\LRfloorDropPct}{3\%}
\newcommand{\RhoLRstanding}{0.30}
\newcommand{\RhoLRcollapse}{0.90}
\newcommand{\RecoveryThirty}{0.334}
\newcommand{\RecoveryFortyFive}{0.549}
\newcommand{\RecoverySixty}{0.700}
\newcommand{\RelLRthirty}{0.112}
\newcommand{\RelLRsixty}{0.490}
\newcommand{\HabermanThirty}{0.35}
\newcommand{\HabermanSixty}{0.74}
\newcommand{\HabermanMarginThirty}{0.05}
\newcommand{\HabCeilThirty}{0.1025}
\newcommand{\HabCeilThirtyCollapse}{0.0241}
\newcommand{\TrueDeltaAssumed}{0.05}
\newcommand{\ObsDeltaThirty}{0.0056}
\newcommand{\NlrThirty}{836}
\newcommand{\NlrSixty}{387}
\newcommand{\LRthirtyVsFloor}{above}
\newcommand{\LRsixtyVsFloor}{above}
\newcommand{\LRthirtyVsFloorHigh}{below}
\newcommand{\HabermanCollapseVerdict}{fails}
\newcommand{\ConfoundSymThirtyFive}{0.59}
\newcommand{\NcontamCheckMid}{290}
\newcommand{\PowerBothContam}{0.87}
\newcommand{\PowerAnywhereContam}{0.81}
\newcommand{\AUCinstr}{0.670}
\newcommand{\AUCinstrLo}{0.651}
\newcommand{\AUCinstrHi}{0.688}
\newcommand{\AUCcomp}{0.636}
\newcommand{\AUCcompLo}{0.617}
\newcommand{\AUCcompHi}{0.656}
\newcommand{\AUCdiff}{0.034}
\newcommand{\AUCdiffLo}{0.020}
\newcommand{\AUCdiffHi}{0.049}
\newcommand{\AUCboth}{0.672}
\newcommand{\PoolBaseRate}{30\%}
\newcommand{\MeasuredHalfwidth}{0.169}
\newcommand{\AnalyticHalfwidth}{0.179}
\newcommand{\HalfwidthRatio}{5.0}
\newcommand{\NsupHi}{520}
\newcommand{\NsupMid}{1555}
\newcommand{\NsupLo}{2590}
\newcommand{\PowerBothSupHi}{0.63}
\newcommand{\PowerCampusSupHi}{0.13}
\newcommand{\NImargin}{0.02}
\newcommand{\NnonInfHi}{207}
\newcommand{\RhoInstr}{0.36}
\newcommand{\RhoComp}{0.29}
\newcommand{\RtwoInstr}{0.130}
\newcommand{\RtwoComp}{0.084}
\newcommand{\RtwoBoth}{0.133}
\newcommand{\DeltaAddInstr}{0.0495}
\newcommand{\DeltaAddComp}{0.0030}
\newcommand{\NaddInstr}{140}
\newcommand{\NaddInstrShift}{222}
\newcommand{\DeltaAddCompVsFloor}{below}
\newcommand{\DichFactor}{0.759}
\newcommand{\DichPenalty}{1.8}
\newcommand{\AUCpenalty}{3.7}
\newcommand{\EffNboth}{157}
\newcommand{\EffNanywhere}{176}
\newcommand{\EffNexpansion}{71}
\newcommand{\EffNthreeCycles}{472}
\newcommand{\EffNscale}{909}
\newcommand{\RoutingDriftPace}{0.098}
\newcommand{\RoutingDriftMAP}{0.000}
\newcommand{\DriftVsEffect}{2.9}
\newcommand{\OfficerMatchPct}{86\%}
\newcommand{\OfficerKappa}{0.43}
\newcommand{\OfficerCeiling}{0.79}
\newcommand{\OfficerKappaShare}{54\%}
\newcommand{\OfficerGateSD}{0.92}
\newcommand{\OfficerImpressionSD}{0.29}
\newcommand{\OfficerRuleHitPct}{40\%}
\newcommand{\OfficerHerHitPct}{33\%}
\newcommand{\OfficerPoolBasePct}{30\%}
\newcommand{\OfficerHitGapPct}{7\%}
\newcommand{\NdisputedPerArm}{742}
\newcommand{\NdecisionsForDispute}{10600}


\title{\textbf{Benchmarking the GT Screener Against CogAT}\\[0.35em]
\large A staged predictive-validity design, its power, and its limits}
\author{GT Selection Capstone --- methodology track}
\date{\today}

\begin{document}
\maketitle

\begin{abstract}
\noindent
GT is building a cognitive screener and needs to know whether it beats CogAT at
the only thing that matters: identifying children who go on to thrive in a
personalised-mastery programme. This report specifies how to answer that
question and, more importantly, reports which parts of the answer are within
reach at the programme's current size and which are not. The criterion is fixed
as mastery pace, anchored to NWEA MAP conditional growth percentiles. The study
is split into a stage that can run now and a stage that waits on CogAT data.
Companion simulation work has since measured the effect sizes involved, and the
result reshapes the recommendation: \textbf{a head-to-head superiority claim
against CogAT is not executable at the current scale} --- it needs roughly
\NsupHi{} to \NsupLo{} children against an effective sample of \EffNboth{} ---
while the better-posed question of whether the screener can \emph{replace}
CogAT is close to reachable now and clearly reachable across two or three
admission cycles. The choice of estimand turns out to matter more than the
sample: framing the comparison as an AUC difference rather than a variance
increment multiplies the required sample by about \AUCpenalty{}. Two assumptions
carry the argument and neither has been measured --- the correlation between
standing ability and learning rate, and the comparator's reliability --- and
both are flagged as contestable rather than settled. Every number here is
generated by a self-testing script in \texttt{scripts/}; none is asserted.
\end{abstract}

\vspace{-0.4em}
\boundary{This is a \emph{predictive-validity} design. It can establish whether
a screener score forecasts how a child does in the programme. It cannot
establish that the programme caused that outcome, and no analysis in this
document should ever be described as measuring GT's impact.
Section~\ref{sec:boundaries} states the line precisely.}

\simnote{Every effect size quoted in this report --- the AUC values, the
variance increments, the routing drift, the officer-agreement figures --- comes
from \textbf{simulated} data produced by a companion synthetic-cohort analysis.
They are design inputs used to size a real study. They are not findings about
real children and must never be quoted as evidence that the screener works.}

\tableofcontents
\clearpage

% ==========================================================================
\part*{Part I --- For decision-makers}
\addcontentsline{toc}{part}{Part I --- For decision-makers}
% ==========================================================================

\section{Executive summary}
\label{sec:exec}

\subsection{What GT can do now}

\keyfinding{\textbf{Start the validity study this school year, against mastery
pace and MAP growth, on the children already enrolled. Do not wait for CogAT.}
This is executable with data GT already collects, it is adequately powered on
the existing \Nbothcohorts{} enrolled children, and it produces a defensible
statement about the screener on its own. It is also time-limited: the open
admission policy that makes it clean will close if GT starts rationing seats
(Section~\ref{sec:range}).}

Three things are executable immediately:

\begin{enumerate}[leftmargin=1.4em,itemsep=0.3em]
\item \textbf{Stage~1 --- does the screener predict who thrives?} On the pooled
cohort, a true validity of $0.35$ is estimated to within about $\pm0.10$
($[\CIbothLo,\ \CIbothHi]$). Section~\ref{sec:stage1}.
\item \textbf{The contamination check.} Whether mastery pace is corrupted by the
platform's own routing is testable now, at power \PowerBothContam{} on the
existing cohort. This matters more than it sounds:
Section~\ref{sec:contamination}.
\item \textbf{The officer-agreement analysis}, the day historical decisions are
released. The record format it needs is specified in Section~\ref{sec:officer}.
\end{enumerate}

\subsection{What GT cannot do now, and what it would take}
\label{sec:exec-constraint}

\keyfinding{\textbf{The head-to-head claim ``our screener beats CogAT'' is not
reachable at the programme's current size, and will not be for two to three
years.} The advantage to be detected is about \AUCdiff{} of AUC. Detecting it
requires \NsupHi{}--\NsupLo{} children depending on how similar the two
instruments are. The effective sample available today, after accounting for
children clustered within campuses and guides, is \EffNboth{}. This is the most
important practical conclusion in the report.}

At the \NcampusCohort{}-child campus cohort the confidence interval on a single
AUC has a half-width of \MeasuredHalfwidth{}, against an effect of \AUCdiff{}
--- the interval is \HalfwidthRatio{} times the effect. The study cannot see its
own effect.

But \NcampusCohort{} is not a fixed ceiling, and the interesting question is
what the addressable sample actually is. Section~\ref{sec:feasibility} works
through every route --- one campus, all campuses after the planned expansion,
the virtual cohort, multiple admission cycles pooled, and the applicant
population at target scale --- and prices each in comparability. The short
answer is that no single-cycle route reaches the superiority requirement, three
pooled cycles reach the variance-increment requirement but not the AUC one, and
only the applicant population at target scale reaches both.

\keyfinding{\textbf{Two changes make the comparison affordable, and neither
involves recruiting more children.} First, ask whether the screener can
\emph{replace} CogAT rather than whether it beats it: a non-inferiority test at
a \NImargin{} AUC margin needs \NnonInfHi{} children rather than \NsupHi{}.
Second, analyse the continuous criterion instead of dichotomising it into
``excels / does not'': the variance-increment test needs \NaddInstrShift{}
children against the artefact floor, roughly \AUCpenalty{} fewer than the AUC
comparison. The estimand is a bigger lever than the sample size, and it is free.}

\subsection{What the simulation found, read honestly}

\simnote{Simulated cohort, not real children.}

The screener reaches an AUC of \AUCinstr{} $[\AUCinstrLo,\ \AUCinstrHi]$ against
mastery pace; the CogAT-like comparator reaches \AUCcomp{}
$[\AUCcompLo,\ \AUCcompHi]$; the paired difference is
$+\AUCdiff{}$ $[+\AUCdiffLo,\ +\AUCdiffHi]$. Both instruments together reach
only \AUCboth{}.

\keyfinding{\textbf{The honest reading is ``not worse, marginally better.''} An
AUC of \AUCinstr{} is a modest instrument. The advantage over the comparator is
real in the simulation but small, and combining the two instruments adds almost
nothing to either. This should not be written up as a win, and any GT-facing
material that describes the screener as outperforming CogAT is ahead of the
evidence --- which at present is simulated evidence in any case.}

There is, however, a genuine asymmetry worth naming. Adding the screener to the
comparator raises explained variance by \DeltaAddInstr{}; adding the comparator
to the screener raises it by \DeltaAddComp{}, which is \DeltaAddCompVsFloor{}
the floor that measurement error alone would produce. In the simulation the
screener largely subsumes the comparator, not the other way round. That is the
strongest defensible form of GT's claim, and it is a \emph{replacement}
argument, not a superiority one.

\subsection{Two assumptions that could overturn this}
\label{sec:exec-assumptions}

A hostile reader will find these, so they are stated here rather than in a
footnote.

\begin{enumerate}[leftmargin=1.4em,itemsep=0.35em]

\item \textbf{The correlation between standing ability and learning rate is
assumed to be \RhoLRstanding{}, and has never been measured by anyone} --- not
by this project, not in the literature. The entire case for the two-phase
instrument depends on it. At the administered block length the learning-rate
readout clears the subscore-reporting bar only while that correlation stays
below \HabermanThirty{}; the working assumption sits just \HabermanMarginThirty{}
below the threshold. At $0.50$ it already fails, and at \RhoLRcollapse{} the
argument collapses at every block length tested. Section~\ref{sec:twophase}.

\item \textbf{The comparator's reliability is assumed to be \RelComparator{}
because this project tried to verify a published figure and could not.} The
primary technical manuals were not text-extractable, so the value was swept
rather than asserted. It matters: the artefact floor that a fake increment would
clear is \SpuriousMid{} at a reliability of \RelComparator{} but \SpuriousWorst{}
at \RelCompLow{} --- a fourfold range driven entirely by an unverified number.
Section~\ref{sec:floor}.

\end{enumerate}

\subsection{The rest, in brief}

\begin{enumerate}[leftmargin=1.4em,itemsep=0.4em]

\item \textbf{The criterion is mastery pace, anchored to MAP --- and the anchor
is not redundant.} In simulation, mastery-pace validity drifted
$+\RoutingDriftPace{}$ AUC purely from changing how the platform routed
children, while MAP conditional growth moved \RoutingDriftMAP{}. That drift is
\DriftVsEffect{} times the entire effect the study is trying to detect. A
criterion that can move by three times the effect size for reasons unrelated to
the child is not a criterion you rely on alone. Section~\ref{sec:criterion}.

\item \textbf{Accuracy is not the binding constraint --- base rate is.} At a
$10\%$ base rate, a screener with validity $0.40$ selecting the top $10\%$ is
right about \PPVFortyTenPct{} of the time. That is arithmetic, and it applies to
CogAT equally. Section~\ref{sec:decision}.

\item \textbf{Not rationing seats is a research asset worth protecting.} A
conventional selective programme sees a validity of $0.40$ attenuated to
\RestrictedTenFromForty{} and needs \NrestrictedForty{} children instead of
\NunrestrictedForty{}. Section~\ref{sec:range}.

\item \textbf{Agreeing with the admissions officer is not the goal, and the
data says so.} A rule using the screener and MAP matches \OfficerMatchPct{} of
her decisions ($\kappa=\OfficerKappa$) --- but a rule fitted to mimic her, using
her own inputs, ceilings at $\kappa=\OfficerCeiling$. No instrument can do
better than that, because the remainder is her own inconsistency. On the files
where the rule and the officer disagree, the rule's picks meet the criterion
\OfficerRuleHitPct{} of the time against her \OfficerHerHitPct{}.
Section~\ref{sec:officer}.

\end{enumerate}

\subsection{Traceability}

Every claim above is checkable. A reader who wants to attack this document
should start at the appendix named beside the claim.

\begin{center}
\small
\begin{tabular}{@{}p{0.58\linewidth}ll@{}}
\toprule
\textbf{Summary claim} & \textbf{Argument} & \textbf{Derivation} \\
\midrule
Superiority against CogAT is not reachable now & \S\ref{sec:feasibility} & App.~\ref{app:auc} \\
Non-inferiority and $\dR$ framings are cheaper & \S\ref{sec:feasibility} & App.~\ref{app:auc} \\
Addressable sample by route & \S\ref{sec:feasibility} & App.~\ref{app:stage1} \\
Mastery pace needs an external anchor & \S\ref{sec:criterion} & App.~\ref{app:contam} \\
Routing contamination is bounded and testable & \S\ref{sec:contamination} & App.~\ref{app:contam} \\
Stage~1 is powered at \Nbothcohorts{}; not at \NcampusCohort{} & \S\ref{sec:stage1} & App.~\ref{app:stage1} \\
CogAT is a concordance question, not equating & \S\ref{sec:stage2} & App.~\ref{app:equating} \\
Unreliability manufactures an increment & \S\ref{sec:floor} & App.~\ref{app:spurious} \\
Base rate, not accuracy, binds the decision & \S\ref{sec:decision} & App.~\ref{app:base} \\
Open admission is worth \NmultiplierTen{}$\times$ the sample & \S\ref{sec:range} & App.~\ref{app:range} \\
Learning rate is not reportable at \BlockAdministered{} trials & \S\ref{sec:twophase} & App.~\ref{app:block} \\
\bottomrule
\end{tabular}
\end{center}

\subsection{What GT must decide, and what GT must supply}

Three decisions are required before the study can be specified further, and none
is a methodological question:

\begin{enumerate}[leftmargin=1.4em,itemsep=0.2em]
\item \textbf{Superiority or replacement?} Whether GT needs to show the screener
\emph{beats} CogAT or only that it can \emph{replace} it. This is the single
largest driver of cost and timeline (Section~\ref{sec:feasibility}).
\item \textbf{What counts as ``excelling''} --- a threshold or a rank. This sets
the base rate, and the base rate determines whether the decision arithmetic in
Section~\ref{sec:decision} is favourable or hostile.
\item \textbf{What the score is for} --- admission, placement, or flagging for
review. A score used to \emph{exclude} must clear a far higher bar than one used
to \emph{surface} children for a human look.
\end{enumerate}

The data required is listed in Section~\ref{sec:data}. Most of it is believed to
exist already; where that belief is unverified it is marked as an assumption
rather than stated as fact.

% ==========================================================================
\section{What this design can and cannot support}
\label{sec:boundaries}

This section exists because the most likely way for this work to fail is not a
statistical error. It is a sentence in a deck that overstates what was found.

\textbf{The design can support:}
\begin{itemize}[leftmargin=1.4em,itemsep=0.15em]
\item that screener scores predict mastery pace and MAP growth, with a stated
confidence interval;
\item that the screener adds (or fails to add) predictive information beyond
CogAT, once CogAT scores exist and once the artefact floor is accounted for ---
subject to the sample being reachable, which Section~\ref{sec:feasibility}
concludes it currently is not for the superiority version of that claim;
\item that decisions made at a given cut have stated sensitivity, specificity
and positive predictive value at a stated base rate;
\item that classification is or is not stable across administrations;
\item that predictive accuracy does or does not differ across demographic
groups.
\end{itemize}

\textbf{The design cannot support:}
\begin{itemize}[leftmargin=1.4em,itemsep=0.15em]
\item that GT \emph{causes} the outcomes observed. Every child in the sample is
enrolled; there is no comparison condition. Predictive validity is a statement
about ordering children, not about the effect of the programme on any of them.
\item that a child who scores highly \emph{would have} done worse elsewhere.
That is a counterfactual, and nothing in this design identifies one.
\item that growth from autumn to spring is attributable to the programme.
Children grow. Conditional growth percentiles control for starting level, not
for maturation, schooling in general, or selection into GT.
\item that a validity coefficient justifies a cut score. Validity is necessary
and not sufficient; the sufficiency argument is the decision arithmetic in
Section~\ref{sec:decision}, and it depends on the base rate and on the
consequences of each error type.
\item that the screener \emph{outperforms} CogAT. At the sample sizes available
this is not a claim the study will be able to make in either direction, and an
inconclusive comparison is not evidence of equivalence unless it was designed as
an equivalence test in advance (Section~\ref{sec:feasibility}).
\end{itemize}

\boundary{Predictive validity is not programme impact. Before-and-after growth
is not a counterfactual. If GT wants an impact claim it needs a different study
with a comparison condition --- a randomised offer, a waitlist, or a
discontinuity --- and none of those is available while the programme admits
everyone. This is a genuine tension: the open-admission policy that makes the
\emph{validity} study unusually clean is exactly what makes an \emph{impact}
study impossible.}

% ==========================================================================
\part*{Part II --- The design}
\addcontentsline{toc}{part}{Part II --- The design}
% ==========================================================================

\section{The criterion problem}
\label{sec:criterion}

Everything downstream depends on this section. A validity coefficient is a
statement about the relationship between two things, and if the second thing is
badly chosen, the coefficient is uninterpretable no matter how large the sample.

\subsection{The ratified choice}

The owner has selected \textbf{mastery pace} --- objectives or lessons mastered
per unit time in the personalised-mastery system --- as the primary criterion,
with \textbf{MAP conditional growth percentile}, autumn to spring, as an
external anchor. This report treats that as settled and spends its effort on
what the choice costs and how to defend it.

\subsection{Why the alternatives were not chosen}

\begin{center}
\footnotesize
\begin{tabular}{@{}p{0.185\linewidth}p{0.115\linewidth}p{0.115\linewidth}p{0.24\linewidth}p{0.215\linewidth}@{}}
\toprule
\textbf{Candidate} & \textbf{Cost to collect} & \textbf{Likely reliability} &
\textbf{Contamination} & \textbf{Gameable?} \\
\midrule
Mastery pace \newline \textit{(primary)} &
Near zero if logged &
High --- many events per child &
\textbf{Serious}: partly a property of platform routing, not the child &
Yes --- by the platform's own tuning, not usually by the child \\
\addlinespace
MAP conditional growth percentile \newline \textit{(anchor)} &
Near zero --- already administered 3$\times$/yr &
High; MAP overall reliability $\approx\RelMAP$ &
Low --- external to routing; conditions on starting RIT &
Low --- third-party normed and proctored \\
\addlinespace
Raw MAP RIT gain &
Near zero &
Lower than CGP --- difference scores &
\textbf{High}: high starters gain fewer points, so it penalises exactly the
children GT selects &
Low \\
\addlinespace
Growth on internal assessments &
Low &
Unknown &
High --- same system that routes the child also assesses them &
Yes \\
\addlinespace
Objectives completed \newline (no time denominator) &
Near zero &
High &
Confounded with enrolment length and attendance &
Yes \\
\addlinespace
Time-to-mastery per objective &
Near zero &
Moderate --- heavy-tailed, needs censoring care &
Same routing problem as pace &
Yes \\
\addlinespace
Persistence / retention &
Near zero &
Low as a per-child measure --- rare, binary, slow &
Dominated by family circumstance, fees, relocation &
No, but uninformative \\
\addlinespace
Teacher or guide rating &
Moderate --- staff time &
Low--moderate; rater effects usually large &
\textbf{Fatal if raters know the screener score} &
Yes \\
\addlinespace
Parent rating &
Moderate &
Low &
Reflects satisfaction more than attainment &
Yes \\
\bottomrule
\end{tabular}
\end{center}

Two entries deserve emphasis because they are the traps.

\textbf{Raw RIT gain is the wrong metric and would understate GT's case.}
NWEA's own documentation notes that high-starting students ``generally do not
raise their score as much,'' and that expected growth compresses in the upper
grades. A gifted cohort scored on raw gain will look mediocre for reasons that
have nothing to do with the programme or the screener. Conditional growth
percentiles exist precisely to solve this, by comparing each child to peers with
the same starting RIT.

\textbf{Teacher rating is the criterion most likely to produce a spuriously
excellent result.} If a guide knows a child was flagged as gifted, the rating
partly measures the flag. Any rating-based criterion must be collected blind to
the screener score, and if that cannot be guaranteed the criterion should be
dropped rather than caveated.

\subsection{Why the anchor is necessary rather than redundant}
\label{sec:anchor}

If mastery pace and MAP growth were near-perfectly correlated, running both
would be waste. They are not, and the anchor earns its place three times over.

\textbf{(a) Conditional growth already solves the starting-level problem.} A
validity study relates a screener score --- which measures standing ability ---
to a growth outcome. If the growth outcome is itself depressed for high
starters, the study builds in a negative bias against precisely the children the
screener identifies, and the resulting validity estimate is too low for a reason
that has nothing to do with the screener. Regression to the mean produces the
same artefact in the opposite direction for low starters. NWEA's conditional
growth percentile is constructed by conditioning on starting RIT, which removes
this confound by design rather than by statistical adjustment after the fact.
Doing the adjustment ourselves would require us to estimate the
growth-on-starting-level relationship from GT's own small sample; NWEA has
estimated it on a normative sample of millions.

\textbf{(b) It survives the ``grading your own homework'' objection.} This is
the audience argument. Mastery pace is a number GT's own platform computes,
using GT's own definition of an objective, for children GT admitted. Every one
of those is a legitimate target for a hostile reviewer, and none of the
rebuttals is fully satisfying. MAP is administered by a third party against
published national norms. A validity coefficient against MAP conditional growth
is not a claim GT can be accused of having constructed.

\textbf{(c) It is not a function of what the platform chose to serve.} This is
the deepest reason and it is the subject of the next subsection.

\subsection{Circularity, and two distinct versions of it}
\label{sec:contamination}

A criterion that is itself a function of the admission decision makes the whole
exercise circular: the study would be measuring how well the screener predicts
its own consequences. There are two ways this can happen here, and they need
different treatment.

\subsubsection*{Circularity of the first kind: the decision changes the treatment}

If admitted children were treated differently \emph{because} of the admission
decision --- placed in a faster track, given more staff attention, held to
different expectations --- then any criterion measured after admission is
contaminated by the decision itself. A screener that predicts admission would
then predict the criterion trivially.

\textbf{The open-admission policy defuses this entirely, and this is an
underrated second benefit of not rationing seats.} If effectively every
applicant is admitted, the screener score does not determine what happens next,
because nothing about placement currently depends on it. The screener can be
scored and sealed, and the programme proceeds as it would have anyway. This is a
clean prospective design, and it is available only while the score is not in
operational use.

\assumption{This holds only if the screener score is genuinely withheld from
guides, from the platform, and from placement decisions during the validation
window. If any part of the programme sees the score, circularity of the first
kind returns. Confirm the score can be sealed operationally.}

\subsubsection*{Circularity of the second kind: the platform routes the child}

This one is live, and it is the reason the anchor exists.

In a personalised-mastery system, how fast a child clears objectives depends on
which objectives they were given. If the platform's routing responds to anything
correlated with the screener signal, then mastery pace partly measures the
routing, and the screener will appear to predict pace for a reason that is not
about the child's ability.

The distinction that matters is \emph{what the routing responds to}:

\begin{itemize}[leftmargin=1.4em,itemsep=0.2em]
\item Routing that responds to the child's \emph{demonstrated performance on the
platform} is benign. That is what a mastery system is supposed to do, and the
resulting pace is a legitimate composite of ability and progress.
\item Routing that responds to an \emph{externally supplied label} --- a prior
CogAT score, a grade placement, a gifted flag, an enrolment tier --- is not
benign. If any such label correlates with the screener, it opens a path from
screener to pace that bypasses the child entirely.
\end{itemize}

\textbf{How much of the primary criterion's validity rests on this
assumption?} It can be bounded. If routing $R$ inflates the observed
screener--pace correlation through a path $\rho_{SR}\cdot\rho_{RP}$, then to
explain away an observed validity of $0.35$ in full, that product must equal
$0.35$ --- meaning both legs would have to be about \ConfoundSymThirtyFive{}. A
routing variable correlating \ConfoundSymThirtyFive{} with a sealed screener
score, and \ConfoundSymThirtyFive{} with pace net of everything else, is a
strong and specific claim. It is not impossible; it is falsifiable.
Appendix~\ref{app:contam} tabulates the full sensitivity surface.

\subsubsection*{How large this effect can be: the simulated demonstration}

\simnote{The result below is a simulation of a mechanism, not a measurement of
GT's platform. Its value is in the comparison of magnitudes, not in the
magnitudes themselves.}

The companion analysis built a cohort, then varied nothing except how the
platform routed children, and re-measured criterion validity under each routing
regime. Mastery-pace validity drifted by $+\RoutingDriftPace{}$ AUC. MAP
conditional growth moved \RoutingDriftMAP{}.

\keyfinding{The routing-induced drift in mastery-pace validity is
\DriftVsEffect{} times the entire effect this study is trying to detect
(Section~\ref{sec:feasibility}). A criterion that can move by three times the
effect size for reasons that have nothing to do with the child is not a
criterion to rely on alone. This is the argument for the external anchor, and it
is quantitative rather than a matter of taste.}

\boundary{\textbf{The flatness of the MAP result is by construction, not a
finding.} MAP was built into the simulation as external to routing, so it could
not move; the \RoutingDriftMAP{} is arithmetic, not evidence. What the
simulation demonstrates is that the \emph{mechanism} is large enough to matter
and that a criterion outside the routing loop is immune \emph{to the extent that
it really is outside the loop}. Whether MAP genuinely is outside GT's routing
loop is an empirical question about GT's platform, answered by the routing-input
audit below and not by any simulation.}

\textbf{The check that would detect it.} Contamination inflates the
screener--pace correlation but cannot inflate the screener--MAP correlation,
because MAP is outside the platform's routing. So the diagnostic is the
\emph{gap} between the two:

\[
H_0:\ \rho(\text{screener},\text{pace}) = \rho(\text{screener},\text{MAP CGP}).
\]

A materially larger correlation with pace than with MAP is the signature of
routing contamination. This is a test of two dependent correlations sharing a
variable, and it is powered on the cohort GT already has: detecting a gap
between $0.45$ and $0.30$ needs \NcontamCheckMid{} children, and the pooled
\Nbothcohorts{} reaches power \PowerBothContam{} (Appendix~\ref{app:contam}).

Three supporting checks, in decreasing order of importance:

\begin{enumerate}[leftmargin=1.4em,itemsep=0.2em]
\item \textbf{Routing-input audit.} Ask the platform team, in writing, exactly
which inputs the routing algorithm consumes. If no externally supplied ability
label is among them, circularity of the second kind is largely closed by
construction, and the statistical checks become confirmatory rather than
load-bearing. This costs one conversation and is worth more than any amount of
modelling.
\item \textbf{Exposure-adjusted validity.} Re-estimate the screener--pace
relationship controlling for routing exposure (objectives served, mean
difficulty served, minutes active, queue length). A large drop in the
coefficient indicates the relationship runs through routing.
\item \textbf{Routing-predictability test.} Regress a routing summary on the
sealed screener score. Some association is expected and benign; the question is
whether it survives conditioning on the child's own on-platform performance.
\end{enumerate}

\assumption{All three supporting checks require routing/exposure telemetry at
child level. Whether the platform retains this is currently unconfirmed. If it
does not, the gap test against MAP remains available and is the one that
matters most.}

\subsection{Measurement detail: do not regress on the raw percentile}

Conditional growth is published as a percentile. Percentile ranks are
approximately uniform by construction, and a uniform variable cannot correlate
perfectly with a normally distributed predictor even under a perfect monotone
relationship: the ceiling is $\sqrt{3/\pi}=\UniformAttenR$. Using raw
percentiles therefore discards about \UniformAttenLossPct{} of the explainable
variance before the analysis begins. Apply a normal-quantile transform to the
conditional growth percentile and analyse on that scale.
Appendix~\ref{app:uniform} derives the constant.

% ==========================================================================
\section{Stage 1 --- validity with no comparator}
\label{sec:stage1}

Stage~1 asks: \emph{does the screener score predict mastery pace and MAP
conditional growth?} It requires no CogAT data and can begin as soon as outcome
data is joined.

\subsection{Design}

Score each enrolled child on the screener. Seal the score. At the end of the
window, join to mastery pace over the same period and to MAP conditional growth
autumn-to-spring. Estimate two correlations and their confidence intervals, and
the gap between them (Section~\ref{sec:contamination}).

Because two primary criteria are tested, control the family-wise error rate;
the sample-size tables report both the unadjusted and the Bonferroni-adjusted
requirement.

\subsection{What the existing cohorts can and cannot deliver}

\begin{center}
\small
\begin{tabular}{@{}rccccc@{}}
\toprule
$n$ & \textbf{Smallest} & \textbf{95\% CI if} & \textbf{CI half-} &
\textbf{Power at} & \textbf{Power at} \\
 & \textbf{detectable $r$} & \textbf{true $r=0.35$} & \textbf{width} &
 $r=0.35$ & $r=0.25$ \\
\midrule
\tabinput{generated/tab_stage1_precision.tex}
\bottomrule
\end{tabular}
\end{center}

\keyfinding{\textbf{The \NcampusCohort{}-child campus cohort is a pilot, not a
study.} Its smallest detectable correlation is \MDEcampus{}. If the true
validity is $0.35$, the interval it returns is
$[\CIcampusLo,\ \CIcampusHi]$ --- an interval that contains ``almost useless''
and ``very good'' simultaneously. Dummy-testing those children is still the
right first move, because it de-risks administration, timing and scoring. It
should not be presented as evidence of validity in either direction.}

The pooled \Nbothcohorts{} children are a different matter: the smallest
detectable correlation falls to \MDEboth{} and the interval narrows to
$[\CIbothLo,\ \CIbothHi]$, which is precise enough to act on. Reaching a
half-width of $\pm0.10$ requires \NciTenAtThirtyFive{} children; $\pm0.075$
requires \NciSevenFiveAtThirtyFive{}.

\subsection{Plan for precision, not for significance}

\begin{center}
\small
\begin{tabular}{@{}cccc@{}}
\toprule
\textbf{True $r$} & \textbf{$n$ for 80\% power} & \textbf{$n$, Bonferroni} &
\textbf{$n$ for CI $\pm0.10$} \\
\midrule
\tabinput{generated/tab_stage1_samplesize.tex}
\bottomrule
\end{tabular}
\end{center}

The two right-hand columns differ sharply, and the difference is the point.
Detecting that validity is non-zero at $r=0.35$ takes \NforThirtyFive{}
children. Estimating it well enough to choose a cut score takes
\NciTenAtThirtyFive{}. Only the second is decision-relevant: nobody will set an
admissions policy on the finding that validity is ``not exactly zero.''

\subsection{Clustering}

Children are nested in campuses and guides, who share practices. Ignoring this
inflates apparent precision. With an intraclass correlation of $0.05$ and all
\NcampusCohort{} campus children under one roof, the design effect is
\DEcampus{} and the effective sample is \EffNcampus{}.

\begin{center}
\small
\begin{tabular}{@{}cccc@{}}
\toprule
\textbf{Cluster size} & \textbf{ICC} & \textbf{Design effect} &
\textbf{Effective $n$ from \Nbothcohorts{}} \\
\midrule
\tabinput{generated/tab_design_effect.tex}
\bottomrule
\end{tabular}
\end{center}

The practical implication favours GT Anywhere: a virtual cohort spread across
many guides clusters far less than a single physical campus. Record a guide or
cohort identifier for every child so this can be estimated rather than assumed.

\subsection{Officer agreement --- specified, pending data}
\label{sec:officer}

Historical admission decisions exist but have not been released. The analysis is
specified here so it can run the day they arrive.

\textbf{Required record format}, one row per historical applicant:
applicant identifier (joinable to the screener); decision date; decision
(admit / decline / waitlist / withdrawn); reviewing officer identifier;
any rubric subscores recorded at the time; and the materials the officer saw.
Cases reviewed by two or more officers must be identifiable as such, because
they are the only rows that estimate inter-officer reliability.

\textbf{Analysis.} Quadratic-weighted $\kappa$ between the screener's implied
decision and the officer's actual decision; inter-officer $\kappa$ on
double-reviewed cases; and the association between officer decisions and the
eventual criterion.

\boundary{Agreement with officers is not validity. If officers' decisions have
little predictive value, a screener that agrees with them perfectly also has
little predictive value --- high agreement would then be evidence \emph{against}
the screener, not for it. Officer agreement is a deployment-risk measure: it
predicts how much disruption adopting the screener would cause. Inter-officer
$\kappa$ on double-reviewed cases is the more useful quantity, because it bounds
how well \emph{any} instrument could agree with a noisy human process.}

\subsubsection*{What the simulation suggests this will look like}

\simnote{Simulated decisions, simulated officer. The pattern below is what the
analysis is expected to \emph{surface}; the numbers are not claims about GT's
actual admissions process.}

A decision rule using the screener and MAP matched \OfficerMatchPct{} of the
simulated officer's decisions, $\kappa=\OfficerKappa$. Read alone that looks
mediocre. It is not the right comparison.

\keyfinding{A rule fitted purely to \emph{mimic} the officer, given her own
stated inputs, ceilings at $\kappa=\OfficerCeiling$. That is the upper bound on
what any instrument could achieve, and the gap below $1.0$ is the officer's own
inconsistency plus weight she places on things outside her stated rule. The
screener-plus-MAP rule therefore reaches \OfficerKappaShare{} of the achievable
agreement, not \OfficerKappa{} of a perfect standard. \textbf{Report agreement
against the ceiling, never against $1.0$} --- doing otherwise holds the
instrument to a standard the human process does not meet.}

The disagreements are more informative than the agreement rate, because they are
structured rather than random, and the two directions mean different things:

\begin{description}[leftmargin=0em,itemsep=0.35em]

\item[Files the rule would admit and she declined] sit \OfficerGateSD{} SD lower
on MAP reading. This is a \emph{documented hard gate} that the statistical rule
does not implement. It is not evidence that the rule is wrong; it is evidence
that a declared policy constraint was left out of the model. The fix is to
implement declared gates explicitly as gates rather than hoping a fitted rule
recovers them --- a hard threshold and a regression weight are different objects,
and a regression will trade a low reading score against a high reasoning score
in exactly the way the gate exists to forbid.

\item[Files she admitted and the rule declined] sit \OfficerImpressionSD{} SD
higher on a construct-irrelevant impression once ability and MAP are partialled
out. This is the finding that matters, and it is the one to handle carefully:
it says some of the variance in her decisions is attributable to something that
is not the construct the process is meant to select on. Whether that is
legitimate expert judgement capturing something the measures miss, or
construct-irrelevant variance in the sense of Messick (1995), cannot be
settled by the agreement analysis. It is settled by asking whether that
component predicts the criterion. If it does not, it is bias.

\end{description}

On the files where the two disagree, the rule's picks reached the criterion
\OfficerRuleHitPct{} of the time against the officer's \OfficerHerHitPct{},
with a pool base rate of \OfficerPoolBasePct{}.

\boundary{\textbf{Do not report that gap as a result.} A difference of
\OfficerHitGapPct{} between two proportions near a third needs about
\NdisputedPerArm{} disputed files \emph{in each direction} to detect at $80\%$
power. Disagreements are roughly one decision in seven, so that is on the order
of \NdecisionsForDispute{} historical decisions --- far beyond what GT will have.
The direction is suggestive and worth pre-registering as a hypothesis. It is not
evidence, and at the sample sizes available it will not become evidence.}

\figslot{officer-agreement}{fig:officer-agreement}{Agreement between screener-implied
decisions and historical officer decisions, with the inter-officer agreement
ceiling shown for reference. The ceiling matters more than the agreement itself:
no instrument can agree with a human process more consistently than that process
agrees with itself. Pending release of historical admission records
(Section~\ref{sec:officer}).}

\figslot{cluster-projection}{fig:cluster-projection}{Low-dimensional projection of
simulated applicant profiles, showing the structure the screener is intended to
resolve and the degree to which candidate criterion groups separate in that
space. Illustrative of the design's logic on synthetic data; it carries no
validity claim about real children.}

\subsection{What Stage~1 licenses on its own}

\boundary{After Stage~1, GT may say: ``the screener predicts mastery pace and
independently-measured academic growth, with validity estimated at $r$ and a
stated interval.'' GT may \emph{not} say that the screener is better than CogAT,
because CogAT was not measured; nor that the screener caused anything; nor that
children scoring highly benefited from the programme. A Stage~1 result is a
statement about the screener, not a comparative or causal one.}

% ==========================================================================
\section{Can the comparison be run at all?}
\label{sec:feasibility}

Everything in Section~\ref{sec:stage2} assumes the CogAT comparison is worth
specifying. Companion simulation work has since measured the effect involved,
and it is small enough that the question of whether GT can afford to detect it
has to be settled before the design is settled. This section does that. It is
placed here, before the Stage~2 design rather than in an appendix, because for a
reader deciding whether to fund this work it is the most consequential section
in the report.

\simnote{The effect sizes below are simulated. Their \emph{magnitude} is
provisional; the \emph{arithmetic} relating an effect size to a required sample
is not, and would apply to any effect of similar size.}

\subsection{The constraint, stated plainly}

The advantage to be detected is a difference of \AUCdiff{} in AUC
($[+\AUCdiffLo,\ +\AUCdiffHi]$) between the screener and a CogAT-like
comparator. Subsampling the simulated cohort down to the size of one physical
campus, \NcampusCohort{} children, gives a $95\%$ interval on a \emph{single}
AUC with a half-width of \MeasuredHalfwidth{}. The interval is
\HalfwidthRatio{} times the effect.

The analytic form corroborates the measurement. Hanley and McNeil's variance for
an AUC (Hanley \& McNeil, 1982) gives a half-width of \AnalyticHalfwidth{} at
\NcampusCohort{} children and a \PoolBaseRate{} base rate --- within $0.01$ of
the subsampled value, which is reassuring about both. Table~\ref{tab:aucprec}
extends it.

\begin{table}[htbp]
\centering
\caption{Precision on a single AUC of \AUCinstr{} at a \PoolBaseRate{} base
rate. The effect to be detected is \AUCdiff{}. Derived in
Appendix~\ref{app:auc}.}
\label{tab:aucprec}
\small
\begin{tabular}{@{}rrrr@{}}
\toprule
$n$ & 95\% half-width & vs.\ the effect & Full interval width \\
\midrule
\tabinput{generated/tab_auc_precision.tex}
\bottomrule
\end{tabular}
\end{table}

\subsection{A correction before despairing}

That comparison is, however, the wrong one, and stating it that way overstates
the difficulty. The half-width of \MeasuredHalfwidth{} is the interval on
\emph{one} AUC estimated on its own. The quantity of interest is the
\emph{difference} between two AUCs measured on the same children, and paired
differences are far more precise than the pair of marginal intervals suggests:
\[
  \operatorname{SE}(A_1-A_2)
  = \sqrt{\operatorname{SE}_1^2 + \operatorname{SE}_2^2
          - 2\,r_{\mathrm{AUC}}\operatorname{SE}_1\operatorname{SE}_2},
\]
where $r_{\mathrm{AUC}}$ is the correlation the two estimates inherit from the
correlation between the two instruments (Hanley \& McNeil, 1983). Two instruments that
both measure reasoning are heavily correlated, and that correlation works in the
study's favour --- it cancels out of the difference.

This matters by a factor of several, so it is worth being explicit: reading two
overlapping confidence intervals and concluding ``no difference'' is a standard
statistical error (Gelman \& Stern, 2006), and the correct paired test is
substantially more powerful. It is not, unfortunately, powerful enough.

\begin{table}[htbp]
\centering
\caption{Sample required for the paired AUC comparison, by how strongly the two
instruments correlate. Power $0.80$, $\alpha=0.05$ two-sided.}
\label{tab:pairedauc}
\small
\begin{tabular}{@{}rrrr@{}}
\toprule
$r_{\mathrm{AUC}}$ & $n$ for power $0.80$ & Power at \NcampusCohort{} &
Power at \Nbothcohorts{} \\
\midrule
\tabinput{generated/tab_paired_auc.tex}
\bottomrule
\end{tabular}
\end{table}

Even under the most favourable assumption --- two instruments so alike that
their AUC estimates correlate at $0.9$ --- the paired test needs \NsupHi{}
children, and reaches power \PowerCampusSupHi{} on one campus and
\PowerBothSupHi{} on every enrolled child GT currently has. If the instruments
are merely moderately similar, the requirement is \NsupMid{} to \NsupLo{}.

\subsection{What the addressable sample actually is}

\NcampusCohort{} is a cohort, not a ceiling, and the recommendation changes
depending on which population the study draws from. Table~\ref{tab:routes}
prices each route.

The raw head count is not the usable sample. Children are nested within campuses
and guides, and children who share a learning environment have correlated
outcomes. The usable sample is $n_{\mathrm{eff}} = n / [1 + (m-1)\rho_I]$, with
$m$ the cluster size and $\rho_I$ the intraclass correlation
(Killip et al., 2004). Pooling several campuses buys fewer independent
observations than it appears to, and pooling one large campus buys very few.

\assumption{An intraclass correlation of $0.05$ is assumed throughout
Table~\ref{tab:routes}; it is a common value in school-effects work but has not
been estimated for this programme. It should be estimated from the first
cohort's data, because the routes differ in how much it costs them. The
expansion route is the most exposed: its design effect is the largest in the
table.}

\begin{table}[htbp]
\centering
\caption{Addressable sample by route. ``Increment'' asks whether the route
reaches \NaddInstrShift{} effective children, enough for the variance-increment
test against the artefact floor; ``AUC'' asks whether it reaches \NsupHi{},
enough for the paired AUC comparison at $r_{\mathrm{AUC}}=0.9$. Enrolment
figures are company-reported (E-101).}
\label{tab:routes}
\small
\begin{tabular}{@{}p{0.34\linewidth}rrrcc@{}}
\toprule
Route & Raw $n$ & Design eff. & $n_{\mathrm{eff}}$ & Increment & AUC \\
\midrule
\tabinput{generated/tab_routes.tex}
\bottomrule
\end{tabular}
\end{table}

Each route costs something different:

\begin{description}[leftmargin=0em,itemsep=0.35em]

\item[One campus, one cycle (\NcampusCohort{}).] Maximum comparability --- one
site, one cohort, one version of the instrument, one set of guides. Minimum
sample. \emph{Effective $n$: \EffNcampus{}.} This route cannot support any
comparative claim.

\item[All physical campuses after the planned expansion ($\approx246$).] GT
intends to open several further campuses. Pooling them multiplies the head count
but introduces the largest design effect in the table, because campuses are big
clusters, and new campuses differ systematically from established ones in staff
experience and cohort composition. \emph{Effective $n$: \EffNexpansion{}} ---
fewer usable observations than the virtual cohort alone, from five times the
children. This route is the trap in the table: it looks like the obvious way to
grow the sample and it is close to the worst.

\item[The virtual cohort, one cycle (\NanywhereCohort{}).] The largest single
clean population, with small clusters and therefore a modest design effect
(\emph{effective $n$: \EffNanywhere{}}), and --- decisively --- it is the
unselective one, which is where the range-restriction advantage of
Section~\ref{sec:range} lives. It costs comparability against the physical
campuses: delivery differs, the criterion is generated by the same platform but
under different supervision, and a claim built here does not transfer to campus
admissions without an invariance argument
(Section~\ref{sec:stage2}).

\item[Every enrolled child, one cycle (\Nbothcohorts{}).] \emph{Effective $n$:
\EffNboth{}.} Requires establishing that the criterion means the same thing in
both settings before pooling --- a measurement-invariance question, not a
convenience. This is the right Stage~1 population and it is still short of the
Stage~2 requirement.

\item[Three consecutive admission cycles ($\approx1{,}038$).] \emph{Effective
$n$: \EffNthreeCycles{}} --- the first route that clears the variance-increment
requirement. It costs time and it costs stability: the instrument must not
change materially across the pooled cycles, the criterion definition must not
change, and NWEA norm updates must be handled. Children who appear in more than
one cycle contribute correlated observations and must be entered once, or
modelled as repeated measures, not counted twice.

\item[The applicant population at target scale ($\approx2{,}000$).]
\emph{Effective $n$: \EffNscale{}} --- the only route that clears both
requirements. It is also the only one where the binding
cost is not enrolment: every child in a Stage~2 analysis must be
\emph{administered CogAT}, and at scale that is a procurement decision rather
than a research one. GT should price it directly --- the required $n$ multiplied
by the per-child cost of a CogAT administration is the actual budget for the
head-to-head claim.

\end{description}

\subsection{The estimand is a bigger lever than the sample}
\label{sec:estimand}

Before concluding that the comparison is unaffordable, it is worth noticing how
much of the requirement is a consequence of how the question was posed rather
than how big the effect is. Table~\ref{tab:framecost} holds the simulated effect
fixed and varies only the estimand.

\begin{table}[htbp]
\centering
\caption{The same simulated effect, four ways of asking about it. Power $0.80$,
$\alpha=0.05$. Derived in Appendix~\ref{app:auc}.}
\label{tab:framecost}
\small
\begin{tabular}{@{}p{0.52\linewidth}rr@{}}
\toprule
Estimand & Effect & $n$ required \\
\midrule
\tabinput{generated/tab_frame_cost.tex}
\bottomrule
\end{tabular}
\end{table}

Two decisions inside that table are free and together move the requirement by
about a factor of \AUCpenalty{}.

\paragraph{Do not dichotomise the criterion.} Mastery pace is continuous and
was made continuous deliberately (Section~\ref{sec:criterion}). Splitting it
into ``excels / does not'' at the \PoolBaseRate{} mark, which is what computing
an AUC requires, multiplies the correlation by \DichFactor{} and therefore
discards about $42\%$ of the explained variance (MacCallum et al., 2002). It costs
a factor of \DichPenalty{} in sample. AUC is a good way to \emph{communicate} an
effect to a non-technical audience and a poor way to \emph{estimate} one; report
it, but do not test on it.

\paragraph{Ask the question GT actually needs answered.} ``Does the screener
beat CogAT?'' and ``can the screener replace CogAT?'' are different questions
with different costs. The second is a non-inferiority question: it asks whether
the screener is not worse by more than some margin GT is willing to accept.

\begin{table}[htbp]
\centering
\caption{Non-inferiority at a \NImargin{} AUC margin against superiority, same
simulated effect, power $0.80$.}
\label{tab:noninf}
\small
\begin{tabular}{@{}rrr@{}}
\toprule
$r_{\mathrm{AUC}}$ & $n$ for non-inferiority & $n$ for superiority \\
\midrule
\tabinput{generated/tab_noninferiority.tex}
\bottomrule
\end{tabular}
\end{table}

At \NnonInfHi{} children the non-inferiority version is within reach of a
pooled cohort today; the superiority version at \NsupHi{} is not. The margin
must be set in advance and justified as a difference GT would genuinely accept,
not chosen after seeing the data --- that is the one place a non-inferiority
design can be abused, and a hostile reader will check it.

\subsection{Verdict}

\keyfinding{\textbf{At the programme's current size, GT can establish that the
screener predicts who thrives, and can plausibly establish that it is not worse
than CogAT. GT cannot establish that it is better.} The superiority claim needs
\NsupHi{}--\NsupLo{} children against an effective sample of \EffNboth{}; it
becomes reachable across three pooled admission cycles for the variance
increment, and only at applicant-population scale for the AUC comparison. This
should shape what GT says publicly about the screener, starting now, rather than
being discovered after the study reads out.}

There is one further reason not to chase the superiority claim, which is that in
the simulation it is not the interesting result. Adding the screener to the
comparator raises explained variance by \DeltaAddInstr{}; adding the comparator
to the screener raises it by \DeltaAddComp{} --- \DeltaAddCompVsFloor{} the
floor that measurement error alone manufactures (Section~\ref{sec:floor}). If
that asymmetry survives contact with real data, the defensible claim is that
CogAT contributes nothing the screener has not already captured. That is a
replacement argument, it is worth more operationally than a marginal AUC
advantage, and it is the cheaper one to establish.

\section{Stage 2 --- incremental validity over CogAT}
\label{sec:stage2}

Stage~2 begins when CogAT scores exist for a substantial subset of the validated
cohort.

\subsection{Why this is not an equating problem}

There is a standing temptation to express screener scores ``in CogAT units.''
That would be equating, and it is not available here. Equating requires that two
tests measure the same construct with comparable reliability and be
interchangeable for the intended purpose; the resulting conversion must hold
regardless of which population is examined. The GT screener and CogAT are
adaptive versus fixed-form, differ in length and reliability, and are being
designed to measure partly different things --- which is the entire premise of
building one.

What \emph{is} available is concordance: a population-dependent statistical
correspondence, honestly labelled as such. But concordance answers the wrong
question. The question GT needs answered is not ``what CogAT score is this
equivalent to,'' it is ``does this tell us anything CogAT does not already tell
us,'' and that is an incremental-validity question.
Appendix~\ref{app:equating} sets out the distinction and its consequences.

\subsection{The model}

Hierarchical regression on the criterion $Y$:

\begin{align}
\text{Step 1:}\quad & Y = \beta_0 + \beta_1 \text{CogAT} + \varepsilon \\
\text{Step 2:}\quad & Y = \beta_0 + \beta_1 \text{CogAT} + \beta_2 \text{Screener} + \varepsilon
\end{align}

The quantity of interest is $\dR$, tested by the $F$-test of the increment. The
effect size to report is $\dR$ itself alongside Cohen's
$f^2=\dR/(1-R^2_{\text{full}})$, plus the semipartial correlation. Report the
increment with a confidence interval, not merely a $p$-value.

Add covariates for grade and age. If both screener components are entered, the
step tests two predictors and the degrees of freedom change accordingly.

\subsection{Sample size}

Section~\ref{sec:feasibility} established the headline: the superiority version
of this comparison is out of reach, and the variance-increment version is the
affordable one. This subsection gives the general table so that the requirement
can be recomputed when GT settles on an increment worth detecting.

\begin{center}
\small
\begin{tabular}{@{}ccccccc@{}}
\toprule
& \multicolumn{3}{c}{\textbf{Screener alone ($q=1$)}} &
  \multicolumn{3}{c}{\textbf{Both components ($q=2$)}} \\
\cmidrule(lr){2-4}\cmidrule(lr){5-7}
$\dR$ & $R^2_{\text{CogAT}}{=}0.10$ & $0.25$ & $0.40$ &
        $0.10$ & $0.25$ & $0.40$ \\
\midrule
\tabinput{generated/tab_stage2_samplesize.tex}
\bottomrule
\end{tabular}
\end{center}

Detecting a three-percentage-point increment over a CogAT baseline of $R^2=0.25$
requires \NincrThreeMid{} children; a five-point increment requires
\NincrFiveMid{}. The simulated increment for adding the screener to the
comparator is \DeltaAddInstr{}, which sits at the friendlier end of that table
and needs \NaddInstr{} children --- roughly what a pooled cohort supplies on a
raw head count, though not after the design effect of
Section~\ref{sec:feasibility} reduces \Nbothcohorts{} heads to \EffNboth{}
usable observations. And that figure holds only \emph{if} the test is run
against a null of zero. It should not be.

\subsection{The artefact floor}
\label{sec:floor}

CogAT is reliable but not perfectly so. Entering an imperfectly measured CogAT
score into step~1 removes only part of the CogAT construct's variance from the
criterion. Whatever remains is available for step~2 to pick up --- and a second
measure of the \emph{same} construct will pick it up.

Appendix~\ref{app:spurious} derives the size of this artefact in closed form.
For two congeneric measures of one latent trait with reliabilities $r_1$ and
$r_2$, where the trait correlates $\rho$ with the criterion and the second
measure contributes \emph{nothing new whatsoever}:

\[
\dR_{\text{spurious}} \;=\; \frac{\rho^{2}\, r_2 \,(1-r_1)^{2}}{1-r_1 r_2}.
\]

\begin{center}
\small
\begin{tabular}{@{}cccc@{}}
\toprule
\textbf{Comparator reliability $r_1$} & $\rho=\RhoCogAToperational$ & $\rho=0.45$ &
$\rho=\RhoCogATobserved$ \\
\midrule
\tabinput{generated/tab_spurious.tex}
\bottomrule
\end{tabular}
\end{center}

\assumption{\textbf{$r_1$ is the second of the two contestable assumptions
flagged in Section~\ref{sec:exec-assumptions}.} A working value of
\RelComparator{} is used here. It is not a verified figure: this project's own
evidence review attempted to confirm a published composite reliability from the
primary technical documentation and could not, because the source PDFs were not
text-extractable. The value was swept forward rather than asserted, and it is
recorded that way in the evidence register. The table is therefore bracketed
from \RelCompLow{} to \RelCompHigh{} rather than quoted at a point, and the
range it produces is not a rounding difference --- the floor moves from
\SpuriousLow{} to \SpuriousWorst{}, a factor of more than sixty across the
corners of the table. \emph{Estimate $r_1$ in GT's own sample rather than
inheriting it.}}

\keyfinding{At the working comparator reliability of \RelComparator{} and a
screener reliability of \RelScreener{}, a screener measuring \emph{nothing}
beyond CogAT still produces an expected increment of \SpuriousMid{} to
\SpuriousHigh{} of criterion variance. At \RelCompLow{} the floor rises to
\SpuriousWorst{}. A study of \NdetectSpurious{} children would detect the
smaller figure at 80\% power and \NdetectSpuriousHigh{} the larger, so at GT's
likely sample the naive test is not \emph{guaranteed} to fire on the artefact
alone --- but the margin between a real increment and a manufactured one is
thin, it depends on a reliability nobody has verified, and it moves in the wrong
direction if the comparator is worse than assumed.}

This is the practical form of a general result: statistically controlling for a
construct by entering one fallible measure of it does not control for the
construct.

\subsection{The fix, and what it costs}

Test against the floor, not against zero:

\[
H_0:\ \dR \le \dR_{\text{spurious}}(r_1, r_2, \rho) \qquad\text{rather than}\qquad H_0:\ \dR = 0.
\]

Equivalently, require the lower bound of the confidence interval on $\dR$ to
exceed the floor. This is more demanding, and the cost is real:

\begin{center}
\small
\begin{tabular}{@{}cccc@{}}
\toprule
\textbf{True $\dR$} & \textbf{$n$, null${=}0$} & \textbf{$n$, null${=}$floor} &
\textbf{Inflation} \\
\midrule
\tabinput{generated/tab_shifted_null.tex}
\bottomrule
\end{tabular}
\end{center}

Detecting a genuine five-point increment against the artefact floor takes
\NshiftedFive{} children rather than \NincrFiveMid{}. A three-point increment
takes \NshiftedThree{}, which is beyond the current cohort. \textbf{GT should
therefore not design Stage~2 to detect small increments.} The defensible target
is a large increment on a moderate sample, not a small increment on a sample
that cannot distinguish it from measurement error.

Three supporting measures, in order of value:

\begin{enumerate}[leftmargin=1.4em,itemsep=0.2em]
\item \textbf{Report the floor alongside the estimate, always.} Whatever else is
done, a $\dR$ reported without its artefact floor is uninterpretable.
\item \textbf{Estimate $r_1$ locally} rather than relying on the publisher's
figure. If CogAT is administered once, use published reliability but treat it as
an assumption and report sensitivity across the
\RelCompLow--\RelCompHigh{} range, as the table above does.
\item \textbf{Prefer a latent-variable model} where sample size allows.
Modelling CogAT and the screener as indicators with fixed reliabilities removes
the artefact by construction rather than by correction. This requires more data
than GT will have at first, so it is the eventual form of the analysis, not the
first one.
\end{enumerate}

\figslot{incremental-validity}{fig:incremental-validity}{Variance in the criterion
explained by CogAT alone and by CogAT plus the screener, with the increment
shown against the artefact floor implied by CogAT's own unreliability
(Section~\ref{sec:floor}). An increment that fails to clear the shaded floor is
consistent with the screener adding nothing. Simulated data illustrating the
comparison the real study will make.}

\subsection{What Stage~2 licenses}

\boundary{After Stage~2, GT may say that the screener does or does not add
predictive information beyond CogAT, by a stated amount, with a stated interval,
against a stated artefact floor. GT may not say the screener is ``better than
CogAT'' without naming the criterion and the population; validity is a property
of a use, not of an instrument. A null result is a real and publishable finding:
it would mean GT can adopt the cheaper, shorter instrument on equivalence
grounds, which is a commercially useful result in its own right.}

\boundary{\textbf{At the sample sizes actually available, Stage~2 will very
likely license the equivalence statement and not the superiority one}
(Section~\ref{sec:feasibility}). This should be settled in advance rather than
discovered at read-out, because the two conclusions want different pre-specified
analyses, and switching to the equivalence framing after a superiority test
fails is exactly the move a hostile reader will look for. Pre-register which
question is being asked, with the non-inferiority margin fixed, before any
CogAT score is linked to a criterion value.}

% ==========================================================================
\section{Decision accuracy, not reliability}
\label{sec:decision}

The screener's output is a decision. The statistics that describe a decision are
not the statistics that describe a score. Coefficient alpha is not among them:
it measures internal consistency, is inflated by test length, and says nothing
about whether the children placed above the cut are the right children.

\subsection{Consistency and accuracy of classification}

Two questions: would a second administration place the child on the same side of
the cut (\emph{decision consistency}), and does the observed classification match
the one the child's true score would produce (\emph{decision accuracy})?

\begin{center}
\small
\begin{tabular}{@{}ccccc@{}}
\toprule
\textbf{Reliability} & \multicolumn{4}{c}{\textbf{Consistency / accuracy at selection ratio}} \\
\cmidrule(lr){2-5}
 & $5\%$ & $10\%$ & $20\%$ & $50\%$ \\
\midrule
\tabinput{generated/tab_decision_consistency.tex}
\bottomrule
\end{tabular}
\end{center}

At a reliability of $0.90$ and a top-$10\%$ cut, consistency is
\DCninetyTen{} and accuracy \DAninetyTen{}. Stated for a leader:
\textbf{about \DisagreeNinetyTenPct{} of children would land on the other side
of the line if tested again next week.} That is not a failure of this
instrument --- it is what a reliability of $0.90$ means, and it applies to CogAT
too. It is an argument for reporting a confidence band with every score and for
routing borderline cases to human review rather than to an automatic decision.

Note that consistency is highest at extreme cuts and lowest at the median
(\DCninetyFifty{} at a $50\%$ cut), because few children sit near an extreme cut.
A high consistency figure at a stringent cut is therefore not as impressive as
it sounds and should not be quoted without the selection ratio.

\subsection{Sensitivity, specificity and the arithmetic that binds}

Once the criterion and the cut are fixed, sensitivity and specificity are not
free parameters --- they follow from the validity, the selection ratio and the
base rate. Modelling the screener and criterion as bivariate normal
(Appendix~\ref{app:base}):

\begin{center}
\small
\begin{tabular}{@{}cccccc@{}}
\toprule
\textbf{Validity} & \textbf{Base rate} & \textbf{Sensitivity} &
\textbf{Specificity} & \textbf{PPV} & \textbf{NPV} \\
\midrule
\tabinput{generated/tab_decision_accuracy.tex}
\bottomrule
\end{tabular}
\end{center}

\keyfinding{At a validity of $0.40$, selecting the top $10\%$, with $10\%$ of
children meeting the ``excels'' criterion: positive predictive value is
\PPVFortyTen{}. \textbf{\FPFortyTenPct{} of the children the screener flags will
not meet the criterion.} The screener is still doing real work --- it
multiplies the hit rate by \PPVliftFortyTen{} over selecting at random --- but
if the score is described
to families as identifying children who will excel, that description is wrong
for most of the children it identifies.}

\subsection{Base rate dominates}

Holding the instrument fixed and moving only the prevalence of ``excelling'':

\begin{center}
\small
\begin{tabular}{@{}cccc@{}}
\toprule
\textbf{Base rate} & \textbf{Se .80 / Sp .90} & \textbf{Se .90 / Sp .95} &
\textbf{Se .95 / Sp .99} \\
\midrule
\tabinput{generated/tab_ppv_baserate.tex}
\bottomrule
\end{tabular}
\end{center}

A very good instrument --- sensitivity $0.90$, specificity $0.95$ --- yields a
PPV of \PPVtwoPercentGoodPct{} when only $2\%$ of children excel,
\PPVfivePercentGoodPct{} at $5\%$, and \PPVtwentyGoodPct{} at $20\%$. The
instrument did not change. Only the definition of ``excels'' changed.

\keyfinding{\textbf{The base rate is a choice GT makes when it defines the
criterion, and it will do more to determine the screener's apparent accuracy
than any amount of psychometric work.} Defining ``excels'' as the top $2\%$
guarantees that most flagged children are false positives, no matter how good
the test is. This should be decided deliberately, in advance, and with the
consequence understood --- not inherited from whatever threshold a report
happens to use.}

The same arithmetic run backwards: to reach a PPV of even $0.50$ at a $5\%$ base
rate requires specificity of \SpecForHalfPPVatFive{}; at a $20\%$ base rate,
\SpecForHalfPPVatTwenty{} suffices.

\begin{center}
\small
\begin{tabular}{@{}ccc@{}}
\toprule
\textbf{Base rate} & \textbf{Specificity for PPV${=}0.50$} &
\textbf{for PPV${=}0.75$} \\
\midrule
\tabinput{generated/tab_required_specificity.tex}
\bottomrule
\end{tabular}
\end{center}

\figslot{decision-accuracy-roc}{fig:decision-accuracy-roc}{ROC curves for the
screener and for CogAT against the ratified criterion, with the proposed
operating cut marked on each. The area under the curve summarises ranking
ability across all cuts; the marked point is what the programme would actually
experience. Simulated data; the real comparison awaits Stage~2.}

\figslot{ppv-base-rate}{fig:ppv-base-rate}{Positive predictive value as a function
of the base rate of ``excelling,'' at fixed sensitivity and specificity. The
curve is steep at low prevalence: the same instrument is trustworthy or
untrustworthy depending on a definitional choice GT has not yet made
(Section~\ref{sec:decision}).}

\subsection{Both error types have owners}

A false positive admits a child who then struggles --- documented in GT's own
history, where children admitted on a non-CogAT ``gifted'' result were later
described as struggling. A false negative denies a child who would have thrived.
These costs fall on different people and are not interchangeable, and the ratio
between them is a values judgement, not a statistical one. The design should
report the full $2\times2$ table at several candidate cuts and let GT choose,
rather than optimising a single figure of merit that hides the trade.

% ==========================================================================
\section{The range-restriction gift}
\label{sec:range}

\subsection{The problem this programme does not have}

Predictive-validity studies in selection are almost always crippled the same
way. Outcomes are observed only for admitted students, and admission depended on
the predictor. The observed correlation is computed in a truncated sample and is
badly attenuated. Corrections exist --- Thorndike's Case~II for direct
restriction on the predictor, Case~III for indirect restriction --- but they
require the unrestricted variance as an input, they assume linearity and
homoscedasticity across a range that was never observed, and they are
contested. Recent work found that standard corrections had systematically
\emph{over}corrected, revising the operational validity of cognitive ability
downward from roughly $0.51$ to $0.31$ --- a revision large enough to overturn
conclusions, produced entirely by the correction machinery.

\textbf{GT currently does not ration seats.} Effectively every applicant is
admitted. Outcomes are observed across the full range of the predictor, and no
correction is needed.

\subsection{What it is worth}

\begin{center}
\small
\begin{tabular}{@{}ccccc@{}}
\toprule
\textbf{Selection} & \textbf{SD ratio} &
\multicolumn{3}{c}{\textbf{Observed $r$ / $n$ needed, by true $\rho$}} \\
\cmidrule(lr){3-5}
\textbf{ratio} & $s/S$ & $\rho=0.30$ & $\rho=0.40$ & $\rho=0.50$ \\
\midrule
\tabinput{generated/tab_range_restriction.tex}
\bottomrule
\end{tabular}
\end{center}

\keyfinding{A true validity of $0.40$ appears as \RestrictedTenFromForty{} in a
sample selected on the top $10\%$ --- an attenuation of
\AttenuationTenPct{}. Detecting it takes \NrestrictedForty{} children instead of
\NunrestrictedForty{}: the open-admission policy is worth a
\NmultiplierTen{}$\times$ reduction in required sample. This is what converts
GT's Stage~1 study from infeasible to comfortably feasible on the cohort that
already exists.}

Even mild selection is expensive: at a top-$20\%$ cut the multiplier is still
\NmultiplierTwenty{}$\times$.

\subsection{How to exploit it, deliberately}

\begin{enumerate}[leftmargin=1.4em,itemsep=0.2em]
\item \textbf{Run the validity study now, while the window is open.} This is
time-limited in a way that is easy to miss. It is the strongest
methodological argument for prioritising Stage~1.
\item \textbf{Record the full predictor distribution}, including children who
enrolled and left, and children who applied and did not enrol. Even under open
admission, self-selection into applying is a form of restriction; documenting
the applicant distribution against a reference population lets its size be
estimated rather than assumed.
\item \textbf{Bank the unrestricted variance now.} The unrestricted standard
deviation of the screener score, measured in an unselected cohort, is the input
every future correction will need. It cannot be recovered later.
\item \textbf{Pre-register the estimate as uncorrected.} The headline validity
should be the raw observed correlation, with no correction applied and none
needed. That is a rare and defensible position and should be stated plainly.
\end{enumerate}

\subsection{If GT becomes selective}

The design degrades in a predictable order. Direct restriction on the screener
becomes Thorndike Case~II, correctable if the unrestricted SD was banked.
Selection on a composite that includes the screener becomes Case~III, which
requires the third-variable correlations as well. Selection on unrecorded
judgement is the bad case: no correction is available, and the validity estimate
becomes a lower bound of unknown tightness.

If selection begins, the mitigations in order of value are: keep scoring every
applicant including those declined; keep the decision rule explicit and
recorded, since an explicit rule is correctable and a judgement call is not; and
if capacity is genuinely oversubscribed, randomise among borderline applicants
--- which would additionally create the only design in this document capable of
supporting a causal claim.

\figslot{range-restriction}{fig:range-restriction}{Attenuation of an observed
validity coefficient as selectivity increases, with GT's current position at the
right-hand edge. The vertical distance between the curves and the horizontal
true-validity line is the bias a conventional study would have to correct away
and GT does not (Section~\ref{sec:range}).}

% ==========================================================================
\section{What the two-phase structure buys}
\label{sec:twophase}

The instrument produces a standing ability estimate and, separately, a
learning-rate estimate from a novel-task block. This section states what that
second component can be claimed to do today, which is nothing, and what it would
take to change that.

\subsection{Current measured state}

Recovery of an injected learning rate on the purpose-built bank, averaged over
eight seeds, is \RecoveryThirty{} at the administered \BlockAdministered{}
trials, rising to \RecoveryFortyFive{} at 45 and \RecoverySixty{} at 60. The
contamination floor --- the apparent learning rate recovered on a scrambled
control arm, where no learning is possible --- is non-zero and, critically,
\emph{does not fall with block length}. The smallest reference spread at which
the readout can name a band is $0.075$ even at 60 trials, against the only
reference distribution with a stated provenance of $0.03$; against that
reference the readout is indeterminate for 100\% of children at every block
length measured.

\boundary{The learning-rate readout is not reportable at the administered block
length. No claim in this report, and none in any GT-facing material, should
assert that the instrument measures learning rate for an individual child. This
is a measured result in GT's own repository, not a cautious hedge.}

\subsection{What it would buy if it reached reportable precision}

The honest case for a learning-rate component is that standing ability measures
what a child has already acquired, and a selection instrument for an
acceleration programme arguably wants to know how fast a child acquires
\emph{new} material. If those are separable, a learning-rate component would
predict criterion variance that standing ability cannot reach, and would be
especially valuable for children whose standing ability is depressed by prior
opportunity rather than by capacity.

That is a hypothesis. The literature on dynamic assessment supports the
construct's coherence but not a specific increment, and the broader finding
across process-signal research is that few studies run the decisive test at all:
the literature is largely \emph{silent} on incremental validity rather than
supportive of it.

\subsection{The right test: when does a subscore earn separate reporting?}

This is exactly the subscore-value question. A subscore should be reported
separately only when it carries information the total score does not, judged by
proportional reduction in mean squared error. And there is a hard ceiling on how
much it can add: the incremental validity of a subscore over the total is no
greater than the reliability of the residual from linear prediction of the
subscore by the total. Appendix~\ref{app:block} derives that ceiling for this
instrument.

\begin{center}
\small
\begin{tabular}{@{}ccccccc@{}}
\toprule
\textbf{Trials} & \textbf{Recovery} & \textbf{Implied} &
\textbf{Max overlap} & \textbf{Ceiling on} & \textbf{Measured} &
\textbf{$n$ to} \\
 & $r$ & \textbf{reliability} & \textbf{with standing} & \textbf{$\dR$} &
 \textbf{$\dR$} & \textbf{detect} \\
\midrule
\tabinput{generated/tab_block_length.tex}
\bottomrule
\end{tabular}
\end{center}

Reading the row that describes the instrument as administered: at
\BlockAdministered{} trials the learning-rate estimate has an implied reliability
of \RelLRthirty{}. It clears the subscore-value bar only if its true correlation
with standing ability is below \HabermanThirty{}.

\subsection{Two claims that must not be conflated}
\label{sec:covariate-vs-absolute}

\simnote{The increments below are simulated.}

The companion analysis measured what the learning-rate readout contributes
against the external MAP anchor: $\dR = \LRdeltaThirty{}$ at the administered
\BlockAdministered{} trials, rising to \LRdeltaSixty{} at 60. It also measured
what lengthening the block does to the two sources of error separately, and the
two behave completely differently.

\begin{center}
\small
\begin{tabular}{@{}lccc@{}}
\toprule
& \textbf{\BlockAdministered{} trials} & \textbf{60 trials} & \textbf{Change} \\
\midrule
Mean posterior SE (random error) & \LRseThirty{} & \LRseSixty{} & $-\LRseDropPct{}$ \\
Contamination floor (systematic) & \LRfloorThirty{} & \LRfloorSixty{} & $-\LRfloorDropPct{}$ \\
\bottomrule
\end{tabular}
\end{center}

\keyfinding{\textbf{Doubling the block halves the random error and does
essentially nothing to the contamination floor.} That single fact separates two
claims which are easy to run together and have entirely different evidentiary
requirements.
\smallskip

\textbf{Claim A --- the learning rate is useful as a covariate.} This asks only
whether the readout \emph{ranks} children in a way that carries criterion
variance. Random error attenuates a ranking but does not bias it, so halving the
posterior SE improves this claim directly, and the measured increment does rise
from \LRdeltaThirty{} to \LRdeltaSixty{}. A longer block buys Claim~A.
\smallskip

\textbf{Claim B --- the instrument measures a child's learning rate.} This is an
\emph{absolute} claim about a quantity attached to a named child, and it is what
anyone reading a score report will assume. It requires the contamination floor to
be small relative to the reported quantity, because a floor is a systematic
offset that does not average out and does not shrink with more trials. The floor
falls by \LRfloorDropPct{} across a doubling of block length. A longer block does
not buy Claim~B, and no feasible block length will.}

Everything that follows about block length is therefore about Claim~A only.
Claim~B is out of reach by a different mechanism than the one a longer block
addresses, and it should be removed from the roadmap rather than deferred to a
larger block.

\subsection{Does the increment survive the artefact floor?}

Even Claim~A is marginal, and whether it survives depends on assumptions rather
than on the measurement.

At \BlockAdministered{} trials the measured increment of \LRdeltaThirty{} sits
\LRthirtyVsFloor{} the \SpuriousMid{} artefact floor that comparator
unreliability manufactures at the working reliability of \RelComparator{}
(Section~\ref{sec:floor}) --- but \LRthirtyVsFloorHigh{} the \SpuriousHigh{}
floor that applies if the criterion correlation is at the upper end of its
plausible range. At 60 trials, \LRdeltaSixty{} sits \LRsixtyVsFloor{} the
working floor. Detecting the \BlockAdministered{}-trial increment against a null
of zero would take \NlrThirty{} children; the 60-trial increment would take
\NlrSixty{}. Neither is available (Section~\ref{sec:feasibility}).

\subsection{The assumption the whole component rests on}

\assumption{\textbf{This is the first of the two contestable assumptions from
Section~\ref{sec:exec-assumptions}, and it is load-bearing for this entire
section.} The correlation between standing ability and learning rate is taken as
\RhoLRstanding{}. It has never been measured --- not in this project, and not,
so far as the evidence review found, anywhere. It is a working value chosen
because it is plausible, not because it is supported.}

The sensitivity is severe, and worse than it first appears.

\begin{center}
\small
\begin{tabular}{@{}ccccc@{}}
\toprule
& \multicolumn{2}{c}{\textbf{\BlockAdministered{} trials}} &
  \multicolumn{2}{c}{\textbf{60 trials}} \\
\cmidrule(lr){2-3}\cmidrule(lr){4-5}
$\rho(\text{standing},\text{rate})$ & Ceiling on $\dR$ & Subscore bar &
Ceiling on $\dR$ & Subscore bar \\
\midrule
\tabinput{generated/tab_lr_sensitivity.tex}
\bottomrule
\end{tabular}
\end{center}

\keyfinding{At the administered block length the subscore-value bar is cleared
only while the correlation stays below \HabermanThirty{}. The working assumption
of \RhoLRstanding{} sits \HabermanMarginThirty{} below that threshold --- and the
component \textbf{already fails at $0.50$}, not at the $0.90$ that would
represent near-redundancy. At \RhoLRcollapse{} it \HabermanCollapseVerdict{} at
every block length in the table, including 60 trials, and the ceiling on its
contribution falls to \HabCeilThirtyCollapse{}.
\smallskip

\noindent A correlation of $0.5$ between two cognitive constructs measured on the
same children by the same instrument is not an extreme scenario; it is an
ordinary one. \textbf{The two-phase argument is one modest empirical result away
from being void, and that result has never been collected.} Measuring this
correlation is the cheapest high-value experiment available to GT and should be
prioritised above lengthening the block, because it determines whether
lengthening the block is worth doing at all.}

\subsection{What follows}

\begin{enumerate}[leftmargin=1.4em,itemsep=0.2em]
\item \textbf{Do not report a learning-rate score to any user.} Not as a number,
not as a band, not as a hedge.
\item \textbf{Do keep administering the block during validation, and keep the
raw trial data.} It costs test time but it is the only way to obtain the two
unknowns that decide the component's fate --- its true correlation with standing
ability, and its true incremental validity --- and those cannot be recovered
retrospectively.
\item \textbf{Treat 60 trials as the minimum plausible reportable length}, and
recognise that even there the contamination floor does not fall. The block
length is necessary, not sufficient.
\item \textbf{Decide the block on test-time economics.} A 60-trial novel block
inside a ${\le}45$-minute assessment competes directly with items measuring
standing ability, which is measured well. The question is not whether learning
rate is interesting; it is whether it is worth more than the standing-ability
precision it displaces.
\end{enumerate}

% ==========================================================================
\section{Minimum additional data}
\label{sec:data}

The owner wants this minimised. The good news is that most of what the design
needs is believed to be collected already. The important caveat is that this
belief is largely unverified.

\assumption{What the programme logs is currently the owner's expectation rather
than a confirmed inventory. Every row below marked \textbf{[A]} is an assumption
to confirm before the study is scheduled. If several prove wrong, the
first-defensible-answer timeline changes materially.}

{\footnotesize
\begin{longtable}{@{}p{0.03\linewidth}p{0.212\linewidth}p{0.245\linewidth}p{0.152\linewidth}p{0.113\linewidth}p{0.072\linewidth}@{}}
\toprule
\textbf{\#} & \textbf{Data element} & \textbf{Why the design needs it} &
\textbf{Already exists?} & \textbf{Cost if not} & \textbf{Stage} \\
\midrule
\endfirsthead
\toprule
\textbf{\#} & \textbf{Data element} & \textbf{Why the design needs it} &
\textbf{Already exists?} & \textbf{Cost if not} & \textbf{Stage} \\
\midrule
\endhead
1 & Screener score per child, sealed &
The predictor. Must be withheld from staff and platform during the window &
Created by the study & Low & 1 \\
\addlinespace
2 & Mastery-pace log: objectives mastered with timestamps &
The primary criterion & \textbf{[A]} Believed logged & Fatal if absent & 1 \\
\addlinespace
3 & MAP RIT autumn + spring, and conditional growth percentile &
The external anchor; the defence against ``grading own homework'' &
Yes --- administered 3$\times$/yr & Low (export) & 1 \\
\addlinespace
4 & Stable child identifier joining screener, platform and MAP &
Without it nothing can be linked & \textbf{[A]} Likely partial &
Moderate --- may need manual reconciliation & 1 \\
\addlinespace
5 & Enrolment start/end and active-days denominator &
Pace is a rate; an uncorrected numerator confounds pace with time enrolled &
\textbf{[A]} Believed available & Low & 1 \\
\addlinespace
6 & Grade and date of birth &
Conditioning, and age-appropriate norming & Near-certain & Negligible & 1 \\
\addlinespace
7 & Guide / campus / cohort identifier &
Clustering; without it precision is overstated & \textbf{[A]} Likely & Low & 1 \\
\midrule
\multicolumn{6}{@{}l}{\textit{--- everything above is sufficient for a first defensible answer ---}} \\
\midrule
8 & Routing and exposure telemetry: objectives served, difficulty, minutes, queue &
The contamination checks in \S\ref{sec:contamination} &
\textbf{[A]} Uncertain; possibly partial &
Moderate & 1 \\
\addlinespace
9 & Demographics: race, sex, EL status, free-lunch or equivalent &
Differential prediction and fairness analysis; required before operational use &
\textbf{[A]} Partially, via enrolment & Low & 1 \\
\addlinespace
10 & CogAT scores for the validated cohort &
The entire comparator. Does not exist. Needs \NaddInstrShift{} effective
children for the increment, \NsupHi{}+ for superiority (\S\ref{sec:feasibility})
&
\textbf{No} --- must be gathered &
High, and it scales with $n$: licensing, proctoring, staff time & 2 \\
\addlinespace
11 & A repeat CogAT administration on a subsample ($n\approx100$) &
Estimates the comparator's reliability locally instead of inheriting the
unverified \RelComparator{}. That figure sets the artefact floor every Stage~2
conclusion is measured against (\S\ref{sec:floor}) &
\textbf{No} & Moderate --- one extra session for a subsample & 2 \\
\addlinespace
12 & Historical admission decisions with officer identifiers &
Officer-agreement analysis (\S\ref{sec:officer}) &
Exists, \textbf{not yet released} & Low once released & 2 \\
\addlinespace
13 & Double-reviewed historical cases, flagged &
The only rows that estimate inter-officer reliability, which bounds the whole
officer analysis & \textbf{[A]} Unknown & Low & 2 \\
\midrule
\multicolumn{6}{@{}l}{\textit{--- everything above is sufficient for a strong answer ---}} \\
\midrule
14 & Second screener administration on a subsample ($n\approx100$) &
Empirical decision consistency instead of the modelled figures in
\S\ref{sec:decision} & No & Moderate --- one retest session & 2 \\
\addlinespace
15 & Guide ratings collected blind to the screener score &
A fallback criterion, and convergent evidence for pace & No &
Moderate --- staff time & 2 \\
\addlinespace
16 & Multi-year mastery pace and MAP for the same children &
Distinguishes durable prediction from a first-year effect & Accrues with time &
Time only & 3 \\
\bottomrule
\end{longtable}}

\textbf{Top three by necessity.} Rows 1--3. The screener score, the mastery-pace
log, and MAP growth. With those three plus an identifier that joins them, Stage~1
runs. Everything else improves the answer or extends it; only these make an
answer possible at all.

\textbf{Row 11 is small and disproportionately valuable.} It is the only row
that converts a contestable assumption into a measurement. Without it, every
Stage~2 result is reported against a floor inherited from a figure this project
could not verify, and a hostile reader gets to choose which end of the range to
argue from. One repeat administration on a hundred children closes that.

\textbf{The single largest risk in this table is row~4.} A validity study is a
join. If children cannot be matched across the screener, the platform and MAP
with high confidence, sample size is lost precisely where it is scarcest, and
match failure is unlikely to be random --- it will correlate with enrolment
churn, which correlates with the criterion. Confirm the join key before anything
else.

% ==========================================================================
\section{Open questions}
\label{sec:open}

These are recorded rather than answered. Inventing answers to them would be the
main way this document could mislead.

\subsection*{On the criterion}
\begin{enumerate}[leftmargin=1.6em,itemsep=0.15em,start=1]
\item Is ``excels'' a threshold or a rank? This sets the base rate, which
Section~\ref{sec:decision} shows dominates everything else about decision
quality.
\item Over what window is mastery pace computed, and how are partial enrolments,
absences and holidays handled in the denominator?
\item Are objectives comparable in size across subjects and grades? If a
mathematics objective takes twice as long as a reading objective, pace measures
subject mix as much as speed.
\item Which MAP subject is the anchor --- mathematics, reading, or a composite?
This must be pre-declared, or the analysis has an undisclosed multiplicity
problem.
\item Is there a floor or ceiling in the pace measure --- children who master
nothing, or who exhaust available content?
\end{enumerate}

\subsection*{On routing and contamination}
\begin{enumerate}[leftmargin=1.6em,itemsep=0.15em,start=6]
\item \textbf{What inputs does the platform's routing algorithm actually
consume?} The single highest-value question in this list. If no external ability
label is among them, circularity of the second kind is largely closed by
construction.
\item Is routing telemetry retained at child level, and for how long?
\item Can the screener score be genuinely sealed from guides and from the
platform for the duration of the validation window?
\end{enumerate}

\subsection*{On CogAT}
\begin{enumerate}[leftmargin=1.6em,itemsep=0.15em,start=9]
\item Which form and level, administered when, and by whom? Form and level
affect reliability and hence the artefact floor in Section~\ref{sec:floor}.
\item Will CogAT be administered to the whole validated cohort or a subsample?
If a subsample, how selected --- because selecting on anything related to the
criterion reintroduces the range restriction the design otherwise escapes. The
number of administrations is also the budget line that decides whether Stage~2
happens at all (\S\ref{sec:feasibility}).
\item \textbf{What is CogAT's composite reliability in a sample like GT's?} This
report uses \RelComparator{} as a working value because an attempt to verify a
published figure from the primary technical documentation failed. It should be
estimated locally rather than inherited. It moves the artefact floor by a factor
of more than sixty across the corners of the sensitivity table
(\S\ref{sec:floor}).
\item \textbf{A discrepancy to resolve:} the evidence register records a GT
commitment to supply ``every student's CogAT and MAP screeners,'' while the
owner states CogAT data does not yet exist. Both may be true --- a commitment to
gather is not a holding of data --- but the register entry should be corrected
so it is not later read as a claim that the data was already available.
\end{enumerate}

\subsection*{On outcome data}
\begin{enumerate}[leftmargin=1.6em,itemsep=0.15em,start=13]
\item Which of the rows marked \textbf{[A]} in Section~\ref{sec:data} actually
exist, at what granularity, and with what history?
\item Is there historical mastery-pace data for children already enrolled? If
so, a retrospective Stage~1 could run immediately --- at the cost that those
children's screener scores would have to be collected now and related to past
outcomes, which reverses the temporal order and weakens the design.
\end{enumerate}

\subsection*{On admissions history}
\begin{enumerate}[leftmargin=1.6em,itemsep=0.15em,start=15]
\item When will historical decisions be released, and in what format
(\S\ref{sec:officer})?
\item Were any applicants declined? If genuinely none were, the officer-agreement
analysis has no variance to explain and should be dropped rather than run.
\item Is the programme likely to become selective during the validation window?
If so, Section~\ref{sec:range} becomes urgent rather than advisory.
\end{enumerate}

\subsection*{On the instrument}
\begin{enumerate}[leftmargin=1.6em,itemsep=0.15em,start=18]
\item What is the true correlation between learning rate and standing ability?
Unmeasured, and it determines whether the learning-rate block can ever earn
separate reporting (\S\ref{sec:twophase}). The value used in this report's
tables, \RhoLRstanding{}, is an illustrative assumption and nothing more.
\item What is the screener's own test--retest reliability? The
\RelScreener{} used throughout is a design target, not a measurement, and every
artefact-floor figure depends on it.
\end{enumerate}

\subsection*{On scale and design parameters}
\begin{enumerate}[leftmargin=1.6em,itemsep=0.15em,start=20]
\item \textbf{Does GT need a superiority claim or a replacement claim?} This is
a business decision, not a methodological one, and it is the largest single
driver of cost: \NsupHi{}--\NsupLo{} children against \NnonInfHi{}
(\S\ref{sec:feasibility}). It should be answered before anything else in this
list.
\item \textbf{Which population will Stage~2 draw from,} and has the
comparability cost of that route been accepted? Table~\ref{tab:routes} prices
each; the expansion route in particular looks better than it is.
\item \textbf{What is the intraclass correlation of the criterion within
campuses and within guides?} Assumed at $0.05$ throughout. It converts head
count into usable sample and no route's verdict in Table~\ref{tab:routes} is
robust to getting it badly wrong. Estimable from the first cohort at no
additional collection cost.
\item \textbf{How strongly do the screener and CogAT correlate?} It sets
$r_{\mathrm{AUC}}$, which moves the required sample for the paired comparison by
a factor of five (\S\ref{sec:feasibility}). Estimable from the first children who
take both.
\item If GT wants a non-inferiority claim, \textbf{what AUC margin would GT
genuinely accept} in exchange for whatever the screener offers over CogAT? This
must be fixed before data are seen, and justified in terms of what GT gains ---
cost, time, breadth, accessibility --- not chosen to make the result come out.
\end{enumerate}

% ==========================================================================
\clearpage
\part*{Part III --- Technical appendices}
\addcontentsline{toc}{part}{Part III --- Technical appendices}
% ==========================================================================

\appendix

% ==========================================================================
\section{Criterion scale: the cost of regressing on a percentile}
\label{app:uniform}

Conditional growth is published as a percentile rank. Suppose the underlying
growth construct $G$ is normally distributed and the screener relates to it
linearly. The published criterion is $P = \Phi(G)$, which is uniform on $(0,1)$.
Even a perfect monotone relationship between $G$ and $P$ is not a perfect linear
one, so regressing on $P$ loses information.

The size of the loss is the correlation between a standard normal variate and
its own probability integral transform. Integrating by parts with $u=\Phi(z)$
and $dv = z\varphi(z)\,dz$, so $v=-\varphi(z)$:

\[
\mathbb{E}[Z\,\Phi(Z)] = \Big[-\Phi(z)\varphi(z)\Big]_{-\infty}^{\infty}
  + \int_{-\infty}^{\infty}\varphi(z)^2\,dz
  = 0 + \frac{1}{2\pi}\int e^{-z^2}dz = \frac{1}{2\sqrt{\pi}}.
\]

Since $\mathbb{E}[Z]=0$ this is also $\operatorname{Cov}(Z,\Phi(Z))$, and
$\operatorname{Var}(\Phi(Z))=1/12$ because $\Phi(Z)$ is uniform. Hence

\[
\operatorname{corr}(Z,\Phi(Z))
 = \frac{1/(2\sqrt{\pi})}{1\cdot\sqrt{1/12}}
 = \frac{\sqrt{12}}{2\sqrt{\pi}} = \sqrt{\tfrac{3}{\pi}} = \UniformAttenR.
\]

The attainable $\Rsq$ is therefore multiplied by $3/\pi = \UniformAttenRsq$:
about \UniformAttenLossPct{} of explainable variance is discarded by using the
percentile directly. Applying $\Phi^{-1}$ to the conditional growth percentile
before analysis recovers it. The constant is verified numerically against
quadrature in \texttt{selftest.py}.

% ==========================================================================
\section{Stage 1: precision and power for a correlation}
\label{app:stage1}

For a bivariate normal sample, $z=\operatorname{artanh}(r)$ is approximately
normal with standard error $1/\sqrt{n-3}$. A two-sided test of
$H_0\!:\rho=0$ at level $\alpha$ with power $1-\beta$ against a true $\rho$
requires

\[
n \;=\; \left(\frac{z_{1-\alpha/2}+z_{1-\beta}}{\operatorname{artanh}\rho}\right)^{2} + 3 ,
\]

and the confidence interval is obtained by transforming
$\operatorname{artanh}(r) \pm z_{1-\alpha/2}/\sqrt{n-3}$ back through $\tanh$.
The routine reproduces Cohen's tabled value of $n=85$ for $\rho=0.30$ exactly;
at $\rho=0.50$ it returns 30 against Cohen's 28, the normal approximation being
mildly conservative at small $n$, which errs toward larger samples and is
therefore safe for planning.

\textbf{Why the report emphasises interval width over power.} A significance
test answers whether validity is distinguishable from zero. No admissions policy
turns on that. The decision-relevant quantity is the width of the interval,
because the cut score and the expected decision accuracy both depend on the
magnitude. This is why Section~\ref{sec:stage1} reports the sample needed for a
$\pm0.10$ interval (\NciTenAtThirtyFive{} children at $r=0.35$) alongside the
much smaller sample needed for significance (\NforThirtyFive{}).

\textbf{Multiplicity.} Two primary criteria are tested. The tables report the
Bonferroni-adjusted requirement ($\alpha=0.025$) as well as the unadjusted one.
Bonferroni is conservative here because the two criteria are positively
correlated; it is used because the correlation between them is unknown, and a
conservative adjustment is easier to defend than an estimated one.

\textbf{Clustering.} With clusters of size $m$ and intraclass correlation
$\mathrm{ICC}$, the variance of an estimate is inflated by the design effect
$1+(m-1)\mathrm{ICC}$, and the effective sample is $n$ divided by that factor.

% ==========================================================================
\section{Criterion contamination: bounds and the detection test}
\label{app:contam}

\subsection{How large a confounder would have to be}

Let $S$ be the sealed screener score, $P$ mastery pace, and $R$ a routing
variable. If $R$ is the only path from $S$ to $P$ other than the child's
ability, the contribution of that path to the observed correlation is
$\rho_{SR}\rho_{RP}$. To account for an observed $r_{SP}$ in full, the product
of the two legs must equal $r_{SP}$; the symmetric case sets both legs to
$\sqrt{r_{SP}}$.

\begin{center}
\small
\begin{tabular}{@{}cccc@{}}
\toprule
\textbf{Observed $r_{SP}$} & \textbf{Both legs equal} &
\textbf{$\rho_{RP}$ if $\rho_{SR}{=}0.30$} &
\textbf{$\rho_{RP}$ if $\rho_{SR}{=}0.50$} \\
\midrule
\tabinput{generated/tab_contamination_sensitivity.tex}
\bottomrule
\end{tabular}
\end{center}

At an observed validity of $0.35$, a confounder would need both legs at
\ConfoundSymThirtyFive{}, or --- if routing correlates only $0.30$ with the
screener --- a routing-to-pace association above $1$, which is impossible. The
reading is that \emph{weak} routing dependence on the screener cannot explain a
moderate validity; only strong dependence can. That is a specific, checkable
claim rather than an unfalsifiable worry.

This is the logic of sensitivity analysis for unmeasured confounding: rather
than asserting no confounding, state how strong confounding would have to be,
and let the reader judge whether that is plausible.

\subsection{The detection test}

Contamination inflates $\operatorname{corr}(S,P)$ but not
$\operatorname{corr}(S,M)$ where $M$ is MAP conditional growth, because MAP is
outside the platform's routing. The diagnostic is therefore the difference
between two correlations that share the variable $S$, tested by the
Meng--Rosenthal--Rubin statistic:

\[
Z=(z_1-z_2)\sqrt{\frac{n-3}{2(1-r_{PM})\,h}},
\qquad
h=\frac{1-f\,\bar r^{2}}{1-\bar r^{2}},
\quad
f=\min\!\left(\frac{1-r_{PM}}{2(1-\bar r^{2})},\,1\right),
\]

with $z_i=\operatorname{artanh}(r_i)$ and $\bar r^{2}=(r_1^{2}+r_2^{2})/2$.
Required samples:

\begin{center}
\small
\begin{tabular}{@{}ccccc@{}}
\toprule
\textbf{$r$(screener, pace)} & \textbf{$r$(screener, MAP)} &
\multicolumn{3}{c}{\textbf{$n$ for 80\% power, by $r$(pace, MAP)}} \\
\cmidrule(lr){3-5}
 & & $0.30$ & $0.50$ & $0.70$ \\
\midrule
\tabinput{generated/tab_contamination_power.tex}
\bottomrule
\end{tabular}
\end{center}

Note the direction of the dependence on $r_{PM}$: the more strongly pace and MAP
agree, the \emph{easier} the discrepancy is to detect, because the comparison
becomes a within-child contrast. Detecting a gap between $0.45$ and $0.30$ at
$r_{PM}=0.50$ requires \NcontamCheckMid{} children; the pooled cohort of
\Nbothcohorts{} reaches power \PowerBothContam{}. GT Anywhere alone
(\NanywhereCohort{} children) reaches \PowerAnywhereContam{}, so the check does
not depend on pooling the two cohorts --- which matters, because pooling a
virtual and a physical cohort introduces exactly the clustering discussed in
Appendix~\ref{app:stage1}.

\subsection{Interpreting the outcome}

\begin{itemize}[leftmargin=1.4em,itemsep=0.15em]
\item \textbf{No detectable gap.} Pace and MAP agree about the screener.
Contamination is bounded below the detectable difference, and the primary
criterion can be used with its full validity estimate.
\item \textbf{Pace materially exceeds MAP.} Consistent with routing
contamination, and also with pace being a more proximal and less noisy measure
of what the screener predicts. These are not distinguishable by this test alone;
the exposure-adjusted analysis in Section~\ref{sec:contamination} is what
separates them. Until it does, the MAP figure should be treated as the headline.
\item \textbf{MAP exceeds pace.} Not contamination. Most likely the pace measure
is noisy, or its denominator is mis-specified.
\end{itemize}

\boundary{This test bounds contamination; it does not eliminate it. A routing
mechanism that happened to affect pace and MAP equally would be invisible to it.
That is implausible given MAP's externality but is not excluded, and it is the
reason the routing-input audit in Section~\ref{sec:contamination} is listed
ahead of the statistical checks.}

% ==========================================================================
\section{Equating, linking and concordance}
\label{app:equating}

The measurement literature distinguishes several kinds of relationship between
scores on two tests, in decreasing order of strength.

\textbf{Equating} produces interchangeable scores. It requires that the two
tests measure the same construct, at comparable reliability, with the same
intended use, and that the conversion be invariant across subpopulations. Only
then may a score on one be substituted for a score on the other without
qualification.

\textbf{Calibration and projection} are weaker forms, applicable when tests
measure the same construct at different reliabilities or difficulties.

\textbf{Concordance} is a purely statistical correspondence between scores on
tests measuring \emph{different} constructs. It is population-dependent: the
concordance table derived in one group need not hold in another. It supports
statements of the form ``in this population, children at this screener score
tended to obtain that CogAT score,'' and nothing stronger.

\textbf{Why this matters for GT.} A concordance table between the screener and
CogAT would be straightforward to produce once both are administered, and it
would be actively harmful if labelled an equating. It would invite the inference
that the screener can substitute for CogAT in contexts where CogAT is required
--- district gifted identification, for instance --- which the evidence would
not support and which would expose GT to a defensible complaint.

More importantly, concordance answers a question GT does not need answered.
Knowing that a screener score of $X$ corresponds to a CogAT SAS of $Y$ tells GT
nothing about whether the screener should replace CogAT. That decision rests on
which instrument better predicts the criterion, and on whether the screener adds
information CogAT lacks. Both are incremental-validity questions, which is why
Section~\ref{sec:stage2} is framed as it is.

\textbf{The invariance question is separate and still required.} Before
operational use, the screener must be checked for differential item functioning
and for measurement invariance across the groups it will be used on, and for
differential prediction --- whether a given score forecasts the same criterion
value in every group. A test can be free of item bias and still predict
differently across groups; these are distinct properties requiring distinct
evidence.

% ==========================================================================
\section{The increment manufactured by unreliability}
\label{app:spurious}

\subsection{Setup}

Let $T$ be a latent ability, standardised. Let the criterion $Y$ be standardised
with $\operatorname{corr}(T,Y)=\rho$. Let two instruments be congeneric
indicators of $T$ with reliabilities $r_1$ and $r_2$:

\[
X_1=\sqrt{r_1}\,T+\sqrt{1-r_1}\,e_1,
\qquad
X_2=\sqrt{r_2}\,T+\sqrt{1-r_2}\,e_2,
\]

with $e_1,e_2$ independent of each other, of $T$, and of $Y$ given $T$. By
construction $X_2$ carries \emph{no information about $Y$ beyond $T$}: it is a
second thermometer for the same temperature. Write

\[
a=\operatorname{corr}(X_1,Y)=\sqrt{r_1}\,\rho,\quad
b=\operatorname{corr}(X_2,Y)=\sqrt{r_2}\,\rho,\quad
c=\operatorname{corr}(X_1,X_2)=\sqrt{r_1r_2}.
\]

\subsection{Derivation}

$R^2$ from $X_1$ alone is $a^2=r_1\rho^2$. For two predictors,

\[
R^2_{\text{full}}=\frac{a^2+b^2-2abc}{1-c^2}.
\]

Substituting, $a^2+b^2 = \rho^2(r_1+r_2)$ and
$2abc = 2\rho^2\sqrt{r_1r_2}\sqrt{r_1r_2}=2\rho^2 r_1r_2$, so the numerator is
$\rho^2(r_1+r_2-2r_1r_2)$ and

\begin{align}
\dR &= \frac{\rho^2(r_1+r_2-2r_1r_2)}{1-r_1r_2}-r_1\rho^2
     = \rho^2\,\frac{r_1+r_2-2r_1r_2-r_1(1-r_1r_2)}{1-r_1r_2} \notag\\[0.3em]
    &= \rho^2\,\frac{r_2-2r_1r_2+r_1^2r_2}{1-r_1r_2}
     = \boxed{\;\rho^{2}\,\frac{r_2\,(1-r_1)^{2}}{1-r_1r_2}\;}
\end{align}

The closed form is checked in \texttt{selftest.py} against $R^2$ computed
directly from the $3\times3$ correlation matrix, at several parameter
combinations, to machine precision.

\subsection{Reading it}

The expression behaves as it should. It vanishes when $r_1=1$: a perfectly
measured incumbent leaves nothing for a redundant challenger to pick up. It
grows with $\rho^2$, so the artefact is largest exactly when the construct
matters most. It grows with $r_2$, so a \emph{more} reliable redundant
challenger produces a \emph{larger} spurious increment --- which is the
counterintuitive part, and the reason a well-built screener is more likely to
produce a misleading positive than a poorly built one.

Substituting GT's figures --- comparator reliability $\RelCompLow$--$\RelCompHigh$,
a screener reliability target of $\RelScreener$, and a criterion correlation
between the Sackett-revised operational value $\RhoCogAToperational$ and the
observed meta-analytic value $\RhoCogATobserved$ --- gives the table in
Section~\ref{sec:floor}: an artefact between $\SpuriousLow$ and $\SpuriousHigh$,
with $\SpuriousMid$ as the central case.

\subsection{Consequence for the test}

At the central case, the conventional test against $\dR=0$ reaches 80\% power to
declare the \emph{artefact} significant at $n=\NdetectSpurious$; at the
pessimistic case, at $n=\NdetectSpuriousHigh$. Since Stage~2 will plausibly run
at a few hundred children, a naive analysis is at material risk of reporting a
significant increment that is entirely an artefact of CogAT's measurement error.
Testing against the floor instead raises the sample needed for a real
five-point increment from $\NincrFiveMid$ to $\NshiftedFive$.

\assumption{The floor depends on $r_2$, the screener's own reliability, which is
currently a design target of \RelScreener{} rather than a measurement. It also
assumes the two instruments are congeneric indicators of a single latent trait,
which is the worst case for GT --- it is the assumption under which the screener
adds nothing. If the screener genuinely measures a partly distinct construct,
the true floor is lower and the test is conservative.}

% ==========================================================================
\section{Decision accuracy and the base-rate arithmetic}
\label{app:base}

\subsection{The $2\times2$ table is implied, not chosen}

Sensitivity and specificity are often treated as properties of an instrument.
For a continuous score cut at a threshold they are not: they follow from the
validity, the selection ratio and the base rate. Model the screener $X$ and the
criterion $Y$ as standard bivariate normal with correlation $\rho$. Select
children with $X>z_x$ where $z_x=\Phi^{-1}(1-\mathrm{SR})$; define ``excels'' as
$Y>z_y$ where $z_y=\Phi^{-1}(1-\mathrm{BR})$. Then

\[
\mathrm{TP}=\Pr(X>z_x,\,Y>z_y),\quad
\mathrm{FP}=\mathrm{SR}-\mathrm{TP},\quad
\mathrm{FN}=\mathrm{BR}-\mathrm{TP},
\]
\[
\mathrm{TN}=1-\mathrm{SR}-\mathrm{FN},\quad
\text{Se}=\frac{\mathrm{TP}}{\mathrm{BR}},\quad
\text{Sp}=\frac{\mathrm{TN}}{1-\mathrm{BR}},\quad
\text{PPV}=\frac{\mathrm{TP}}{\mathrm{SR}} .
\]

The joint probability is a bivariate normal orthant, computed here by
one-dimensional quadrature and validated against Sheppard's theorem, against
Plackett's identity, and against the published Taylor--Russell tables.

A structural consequence worth noting: sensitivity is bounded above by
$\mathrm{SR}/\mathrm{BR}$. If GT selects the top $10\%$ but $20\%$ of children
meet the criterion, sensitivity cannot exceed $0.50$ however good the instrument
is, because there are not enough places. The low sensitivity figures in
Section~\ref{sec:decision} are largely this, not instrument failure.

\subsection{The odds form, as a cross-check}

The same arithmetic in likelihood-ratio form makes the base-rate dependence
transparent:

\[
\underbrace{\frac{\text{PPV}}{1-\text{PPV}}}_{\text{posterior odds}}
=\underbrace{\frac{\mathrm{BR}}{1-\mathrm{BR}}}_{\text{prior odds}}
\times
\underbrace{\frac{\text{Se}}{1-\text{Sp}}}_{\text{likelihood ratio}} .
\]

The instrument contributes only the likelihood ratio. At Se $=0.90$ and
Sp $=0.95$ the likelihood ratio is $18$, which is strong. At a base rate of
$2\%$ the prior odds are $1{:}49$, so the posterior odds are $18{:}49$ --- a
PPV of \PPVtwoPercentGood{}. \textbf{A likelihood ratio of 18 is an excellent
diagnostic and still leaves most flagged children below the bar}, because the
prior was so unfavourable. No improvement in the instrument fixes this; only
redefining the criterion, or accepting that the score surfaces candidates rather
than identifies successes, does.

\subsection{Decision consistency and accuracy}

Model observed score $X=T+E$ with $\operatorname{Var}(X)=1$ and
$\operatorname{Var}(T)=r_{xx}$. Two parallel administrations correlate $r_{xx}$,
so consistency at a cut $c$ is
$\Pr(X_1>c,X_2>c)+\Pr(X_1<c,X_2<c)$ under a bivariate normal with correlation
$r_{xx}$. Accuracy compares the observed classification to the true-score
classification at the same raw cut, using
$\operatorname{corr}(X,T)=\sqrt{r_{xx}}$ and the cut $c/\sqrt{r_{xx}}$ on the
standardised true score.

At a median cut this has the exact closed form
$\tfrac12+\arcsin(r_{xx})/\pi$, which \texttt{selftest.py} verifies to nine
decimal places. These are model-based figures assuming normality and classical
error; the operational estimate for a real test would use the
Livingston--Lewis procedure on the actual score distribution. The modelled
figures are adequate for design, and the report labels them as modelled.

% ==========================================================================
\section{Range restriction}
\label{app:range}

\subsection{Thorndike Case II}

Under direct selection on the predictor with unrestricted SD $S$ and restricted
SD $s$, write $U=S/s$. The observed correlation in the restricted group relates
to the unrestricted one by

\[
r_{\text{restricted}}
= \frac{\rho/U}{\sqrt{1-\rho^{2}+\rho^{2}/U^{2}}},
\qquad
\rho
= \frac{r_{\text{restricted}}\,U}{\sqrt{1-r_{\text{restricted}}^{2}+r_{\text{restricted}}^{2}U^{2}}} .
\]

For top-fraction selection with ratio $p$ and cut $c=\Phi^{-1}(1-p)$, the
truncated variance follows from the standard normal moments. With
$\lambda=\varphi(c)/p$ the inverse Mills ratio,

\[
\operatorname{Var}(X\mid X>c) = 1+\lambda c-\lambda^{2} .
\]

At $p=0.10$: $c=1.2816$, $\lambda=1.7550$, and the SD ratio is
$s/S=\SDratioTen$. The closed form is checked against numerical quadrature of
the truncated density in \texttt{selftest.py}, and the attenuation and
correction functions are checked to invert one another.

\subsection{What GT avoids}

A true validity of $0.40$ observed under top-$10\%$ selection appears as
$\RestrictedTenFromForty$, an attenuation of \AttenuationTenPct{}. The sample
needed to detect it at 80\% power rises from $\NunrestrictedForty$ to
$\NrestrictedForty$ --- a factor of \NmultiplierTen{}. At top-$20\%$ selection
the factor is still \NmultiplierTwenty{}.

The deeper benefit is that no correction is applied at all. Corrections require
the unrestricted SD as an input, assume linearity and homoscedasticity across a
range that was by definition never observed, and have been shown to
systematically overcorrect: the operational validity of cognitive ability was
revised downward from roughly $0.51$ to $0.31$ once the overcorrection was
identified. \textbf{An uncorrected estimate from an unrestricted sample is not
merely more efficient --- it is a different and stronger class of evidence,
because it does not depend on the correction machinery that produced that
error.}

\subsection{Indirect restriction, if it appears}

If selection operates on a composite that includes the screener rather than on
the screener alone, Case~II no longer applies and Case~III is required, which
additionally needs the correlations between the screener, the composite and the
criterion. Case~III corrections are substantially less stable. The mitigation is
procedural rather than statistical: keep the selection rule explicit and
recorded. An explicit rule is correctable; an unrecorded judgement is not.

% ==========================================================================
\section{Block length and the subscore-value bar}
\label{app:block}

\subsection{From recovery to reliability}

The repository reports \emph{recovery} $r$: the correlation between the
estimated learning rate and the injected true value in simulation. If the
estimate is the truth plus independent error, then
$\operatorname{corr}(\hat\lambda,\lambda)=\sqrt{\text{reliability}}$, so the
implied reliability of the readout is $r^{2}$. At the administered
\BlockAdministered{} trials, recovery of \RecoveryThirty{} implies a reliability
of \RelLRthirty{}; at 60 trials, \RecoverySixty{} implies \RelLRsixty{}.

\subsection{The subscore-value bar}

A subscore earns separate reporting only when it predicts its own true score
better than the total does --- the proportional-reduction-in-mean-squared-error
criterion. Writing $r_s$ for the subscore reliability, $r_x$ for the total's,
and $\rho$ for the correlation between their true scores, the condition is
$r_s > \rho^{2}r_x$, so the subscore clears the bar only when

\[
\rho \;<\; \sqrt{r_s/r_x}.
\]

At \BlockAdministered{} trials this ceiling is \HabermanThirty{}: the
learning-rate block earns separate reporting only if learning rate and standing
ability are correlated below \HabermanThirty{} in truth. At 60 trials the
ceiling relaxes to \HabermanSixty{}.

\subsection{The hard ceiling on incremental validity}

Separately, there is an upper bound on how much a subscore can add to prediction
of an \emph{external} criterion: the incremental validity is no greater than the
reliability of the residual from linear prediction of the subscore by the total.
With $S$ and $X$ standardised, $\operatorname{corr}(S,X)=c=\rho\sqrt{r_sr_x}$,
the residual is $e=S-cX$ with $\operatorname{Var}(e)=1-c^{2}$. Its true part is
$T_e=T_s-cT_x$, and since $\operatorname{Var}(T_s)=r_s$,
$\operatorname{Var}(T_x)=r_x$ and $\operatorname{Cov}(T_s,T_x)=c$,

\[
\operatorname{Var}(T_e)=r_s-2c^{2}+c^{2}r_x=r_s\big[1-\rho^{2}r_x(2-r_x)\big],
\]
\[
\dR_{\max}=\operatorname{rel}(e)
=\frac{r_s\big[1-\rho^{2}r_x(2-r_x)\big]}{1-\rho^{2}r_sr_x} .
\]

At $\rho=0$ this reduces to $r_s$: an orthogonal subscore can add at most its
own reliability, which at \BlockAdministered{} trials is \RelLRthirty{}. At
$\rho=\RhoLRstanding$ the ceiling falls to \HabCeilThirty{}.

\subsection{Sensitivity to the unmeasured correlation}

Both bars --- the subscore-value bar $\rho < \sqrt{r_s/r_x}$ and the incremental
ceiling above --- are governed by $\rho$, the true correlation between learning
rate and standing ability, which has never been measured. The sensitivity is
tabulated in Section~\ref{sec:twophase}. Two features of it deserve emphasis
here, because both are easy to miss from the main text.

First, the failure point at the administered block length is $\rho \approx
\HabermanThirty{}$, and the working assumption of \RhoLRstanding{} sits only
\HabermanMarginThirty{} below it. The component does not survive to
\RhoLRcollapse{} and then fail; it fails at $0.50$, which is an unremarkable
correlation between two cognitive constructs measured on the same children by
the same instrument on the same day.

Second, the two bars fail in different ways and the difference matters for how a
negative result should be reported. Exceeding the subscore-value bar means the
subscore should not be reported \emph{at all}, because the total score predicts
it better than it predicts itself. Falling below the incremental ceiling means
only that the subscore's contribution is bounded and small. The first is a
prohibition; the second is a magnitude. At \RhoLRcollapse{} the administered
block \HabermanCollapseVerdict{} the first, which is the stronger and more
damaging outcome.

\subsection{Attenuation, and the decisive comparison}

The ceiling is an upper bound. The \emph{expected} increment is governed by
attenuation: a predictor with reliability $r_s$ realises only $r_s$ of a true
$\dR$. A genuinely real increment of \TrueDeltaAssumed{} is realised as
$\RelLRthirty\times\TrueDeltaAssumed=\ObsDeltaThirty$ at \BlockAdministered{}
trials.

The companion simulation measured the realised increment directly, against the
external MAP anchor, and obtained \LRdeltaThirty{} at \BlockAdministered{}
trials and \LRdeltaSixty{} at 60. The agreement with the attenuation prediction
at \BlockAdministered{} trials is close, which is mild corroboration that the
attenuation model is describing the right mechanism.

\keyfinding{\textbf{Whether the measured increment clears the artefact floor
depends on an assumption, not on the measurement.} At the working comparator
reliability of \RelComparator{}, the floor is \SpuriousMid{} and the
\BlockAdministered{}-trial increment of \LRdeltaThirty{} sits
\LRthirtyVsFloor{} it. If the criterion correlation is at the upper end of its
plausible range the floor rises to \SpuriousHigh{} and the same increment sits
\LRthirtyVsFloorHigh{} it. If the comparator's reliability is \RelCompLow{}
rather than \RelComparator{}, the floor reaches \SpuriousWorst{} and the
increment is not close.
\smallskip

\noindent The honest summary is that the learning-rate component's contribution
at the administered block length is \emph{the same order of magnitude as the
artefact a second fallible measure of the same construct produces for free}. It
cannot be separated from that artefact without first pinning down the
comparator's reliability in GT's own sample --- which is a measurement GT can
make, and should, before running this comparison at all.}

At 60 trials the measured increment rises to \LRdeltaSixty{}, \LRsixtyVsFloor{}
the working floor, and detectable at \NlrSixty{} children rather than
\NlrThirty{}. Sixty trials is therefore the minimum plausible length for the
covariate claim --- necessary, and still not sufficient for an absolute claim,
because the contamination floor falls only \LRfloorDropPct{} across that
doubling while the random error falls \LRseDropPct{}
(Section~\ref{sec:covariate-vs-absolute}).

\assumption{The correlation between learning rate and standing ability,
$\rho=\RhoLRstanding$ in these tables, is an illustrative value. It has not been
measured and it drives both the subscore-value ceiling and the incremental
ceiling. Measuring it is the main reason to keep administering the block during
validation even though its output is not reportable.}

% ==========================================================================
\section{ROC precision, paired comparison, and the cost of the frame}
\label{app:auc}

This appendix supports Section~\ref{sec:feasibility}. It derives the precision
of an AUC estimate, the precision of the \emph{difference} between two AUCs
measured on the same children, and the translation between the AUC frame and
the variance frame used everywhere else in this report.

\subsection*{Precision of a single AUC}

For an AUC of $A$ estimated on $n_1$ positive and $n_0$ negative cases,
Hanley and McNeil (1982) give
\[
  \operatorname{Var}(A) = \frac{A(1-A) + (n_1-1)(Q_1 - A^2) + (n_0-1)(Q_2 - A^2)}
                               {n_1 n_0},
\]
\[
  Q_1 = \frac{A}{2-A},\qquad Q_2 = \frac{2A^2}{1+A},
\]
where $Q_1$ and $Q_2$ are the probabilities that two positives both outrank a
negative, and that a positive outranks two negatives. The expression is exact
under an exponential model for the score distributions and is used here as an
approximation.

Two checks are run against this in \path{scripts/selftest.py}. At $A=0.5$ it must
reduce exactly to the null variance of the Mann--Whitney statistic,
$\operatorname{Var} = (N+1)/(12\,n_1 n_0)$, which it does to machine precision;
and at $A=1$ it must vanish, which it does. At $n=\NcampusCohort{}$ with a
\PoolBaseRate{} base rate it returns a half-width of \AnalyticHalfwidth{},
against \MeasuredHalfwidth{} measured by subsampling the simulated cohort. The
agreement is close enough to treat either as the working figure.

\subsection*{Why the paired difference is far more precise}

The relevant quantity is not the width of either marginal interval but the
precision of $A_1 - A_2$ when both are computed on the \emph{same} children:
\[
  \operatorname{Var}(A_1 - A_2)
   = \operatorname{Var}(A_1) + \operatorname{Var}(A_2)
     - 2\,r_{\mathrm{AUC}}\sqrt{\operatorname{Var}(A_1)\operatorname{Var}(A_2)},
\]
with $r_{\mathrm{AUC}}$ the correlation between the two AUC statistics, itself
induced by the correlation between the two instruments' scores within the
positive and negative groups (Hanley \& McNeil, 1983). Because both instruments
measure overlapping reasoning constructs, $r_{\mathrm{AUC}}$ is large and the
shared sampling error cancels.

This is the reason Section~\ref{sec:feasibility} does not simply divide the
single-AUC half-width by the effect. Comparing two marginal intervals and
concluding ``no difference'' where they overlap is a well-documented error
(Gelman \& Stern, 2006). The correct paired test is several times more
efficient --- and still, at $r_{\mathrm{AUC}}=0.9$, requires \NsupHi{} children
for $80\%$ power against an effect of \AUCdiff{}.

$r_{\mathrm{AUC}}$ is not known for this pair of instruments and is treated as a
sensitivity parameter throughout, spanning $0.5$ to $0.9$. It can be estimated
directly once any sample has both scores, and doing so early is worthwhile: it
moves the required sample by a factor of five.

\subsection*{Non-inferiority}

A superiority test asks whether the lower bound on $A_1 - A_2$ exceeds zero. A
non-inferiority test asks whether it exceeds $-\delta$ for a pre-specified
margin $\delta$. With a true difference $\Delta$ the required sample follows from
\[
  \frac{\Delta + \delta}{\operatorname{SE}(A_1-A_2)} \;\ge\; z_{1-\alpha} + z_{1-\beta},
\]
one-sided at $\alpha = 0.025$. At $\delta = \NImargin{}$ this gives
\NnonInfHi{} children at $r_{\mathrm{AUC}}=0.9$ against \NsupHi{} for
superiority. The margin must be fixed before the data are seen and justified as
a difference GT would genuinely tolerate in exchange for whatever the screener
offers over CogAT --- cost, time, breadth, or accessibility. A margin chosen
afterwards to make the result come out is the standard abuse of this design and
is easy for a reviewer to detect from the analysis timestamp.

\subsection*{Translating between AUC and $R^2$}

The rest of this report works in correlations and variance increments; the
companion simulation reports AUCs. The bridge, for a continuous predictor $X$
and a latent criterion $Y$ that are bivariate normal with correlation $\rho$,
with the binary outcome defined as $Y$ exceeding the $1-p$ quantile $c$:
\[
  \mathrm{AUC} = \int_{-\infty}^{\infty} f_{+}(x)\,F_{-}(x)\,\mathrm{d}x,
  \qquad
  f_{+}(x) = \frac{\phi(x)\,\bar\Phi\!\big(\tfrac{c-\rho x}{\sqrt{1-\rho^2}}\big)}{p},
\]
and $f_{-}$ the same with $\Phi$ in place of $\bar\Phi$ and $1-p$ in the
denominator. This is evaluated by quadrature in \path{scripts/statlib.py} and
inverted by bisection.

It is validated against an independent reduction. Writing $U = X_i - X_j$, the
event $\{X_i > X_j,\ Y_i > c,\ Y_j \le c\}$ is an orthant probability of the
trivariate normal $(U/\sqrt2,\ Y_i,\ -Y_j)$ with correlations
$(\rho/\sqrt2,\ \rho/\sqrt2,\ 0)$; conditioning on the first coordinate reduces
it to a one-dimensional integral over the bivariate orthant routine already
validated against Sheppard's theorem and Plackett's identity. The two paths
agree to $10^{-6}$.

Applying it to the simulated results: an AUC of \AUCinstr{} at a \PoolBaseRate{}
base rate corresponds to $\rho = \RhoInstr$, and \AUCcomp{} to
$\rho = \RhoComp$. The first of those is worth noting --- the validity of
$0.35$ assumed throughout this report as a working value, chosen before the
simulation was run, lands within $0.01$ of what the simulation produced.

\subsection*{What dichotomising the criterion costs}

Computing an AUC requires a binary criterion. Mastery pace is continuous, so
using AUC means dichotomising it. For a bivariate normal pair split at base
rate $p$, the correlation with the dichotomised criterion is
\[
  r_{\mathrm{pb}} = \rho\,\frac{\phi(c)}{\sqrt{p(1-p)}},
\]
whose multiplier is at most $\sqrt{2/\pi} = 0.798$ at a median split and
\DichFactor{} at $p = \PoolBaseRate{}$. Squaring, about $42\%$ of the explained
variance is discarded --- information the study paid to collect and then threw
away at the analysis step (MacCallum et al., 2002). In sample terms the penalty
is a factor of \DichPenalty{}, and the further step to a paired AUC comparison
costs a factor of \AUCpenalty{} against the continuous variance increment.

The recommendation that follows is narrow and worth stating precisely: compute
AUC, report AUC, and put AUC in the slide deck, because it is the statistic a
non-technical audience can interpret. Do not \emph{test} on it, and do not
\emph{size the study} on it.

\subsection*{The asymmetry between the two increments}

Finally, the two directions of the incremental question are not the same
question and do not have the same answer. In the simulation:
\[
  R^2_{\text{comparator}} = \RtwoComp, \qquad
  R^2_{\text{screener}} = \RtwoInstr, \qquad
  R^2_{\text{both}} = \RtwoBoth.
\]
Adding the screener to the comparator gains \DeltaAddInstr{}; adding the
comparator to the screener gains \DeltaAddComp{}, which falls
\DeltaAddCompVsFloor{} the artefact floor of Appendix~\ref{app:spurious}. A
result of that shape does not license ``the screener is better.'' It licenses
``the comparator adds nothing measurable once the screener is in the model,''
which is a claim about redundancy rather than superiority, and is the claim
Section~\ref{sec:feasibility} recommends GT pursue.

\section{Reproducibility}
\label{app:repro}

\subsection{How the numbers were produced}

Every numeric value in this report is generated by
\texttt{scripts/compute.py} and written into \texttt{generated/} as \LaTeX{}
macros and table bodies, which the document then includes. No number is typed
into the prose by hand. To regenerate:

\begin{center}
\texttt{cd scripts \&\& python3 selftest.py \&\& python3 compute.py}
\end{center}

\texttt{scripts/statlib.py} implements the statistical routines using only the
Python standard library --- no NumPy, no SciPy --- so the results reproduce on a
bare Python 3.8+ interpreter. The routines include the regularised incomplete
beta function, the central and noncentral $F$ distributions, bivariate normal
orthant probabilities, and the power and sample-size functions built on them.

\subsection{Validation}

\texttt{scripts/selftest.py} checks those routines against external reference
values and internal consistency identities, and exits non-zero on any failure.
It is a prerequisite of the build. The external checks include:

\begin{itemize}[leftmargin=1.4em,itemsep=0.1em]
\item central $F$ quantiles against published tables at three
$(\mathrm{df}_1,\mathrm{df}_2,\alpha)$ combinations;
\item two G*Power worked examples for the $\Rsq$-increment test, reproduced
exactly ($n=77$ for $f^2=0.15$ with $q=3$; $n=395$ for $f^2=0.02$ with $q=1$);
\item Cohen's tabled sample size for a correlation ($n=85$ at $\rho=0.30$);
\item Taylor--Russell tabled success ratios at two parameter combinations;
\item Sheppard's theorem for the bivariate normal at the median cut, across five
correlations;
\item Plackett's identity for off-centre orthant probabilities;
\item the Hanley--McNeil AUC variance reduced to the Mann--Whitney null variance
$(N+1)/(12 n_1 n_0)$ at $A=0.5$, exactly, at three sample splits;
\item the AUC/correlation bridge against an independent trivariate-orthant
reduction of the same probability, agreeing to $10^{-6}$
(Appendix~\ref{app:auc}).
\end{itemize}

Internal checks confirm that the range-restriction correction inverts the
attenuation exactly, that the closed-form spurious increment matches $\Rsq$
computed from the correlation matrix, that the $\sqrt{3/\pi}$ constant matches
quadrature, that a zero-validity predictor selects exactly at the base rate, and
that a zero-reliability test agrees with itself only at chance.

\subsection{Building the document}

\texttt{make} regenerates the numbers and builds \texttt{main.pdf}. If no local
\LaTeX{} installation is present, \texttt{make docker} builds the same PDF in a
TeX~Live container. \texttt{make check} runs the self-test and a consistency
check that every macro referenced in the sources is defined in
\texttt{generated/numbers.tex}. That check catches three things the build
either misses or reports late: a stale \texttt{generated/} directory, a macro
that is computed but no longer referenced anywhere (usually a number left behind
by a rewritten passage), and a passage calling a macro that no longer exists.

\subsection{Figures}

The six figures are produced by a companion synthetic-data analysis and are
included via \texttt{\textbackslash figslot}, which falls back to a labelled
placeholder when the PDF is absent. The document therefore builds with or
without them, and no argument in the text depends on a figure's values.

\subsection{Limits of these computations}

Three caveats a reviewer should hold onto. The power calculations assume
bivariate normality and homoscedasticity; mastery pace is a rate and is likely
right-skewed, so the analysis should either transform it or use a model
appropriate to counts, and the sample sizes here should be treated as lower
bounds. The noncentrality convention is $\lambda=f^{2}n$, matching G*Power and
verified against its worked examples; other conventions give slightly different
answers and the difference is not negligible at small $n$. And every figure
depending on the screener's reliability inherits the status of that number,
which is a design target rather than a measurement.

A fourth caveat is more consequential than the other three and is stated in the
main text as well. The two parameters that most strongly drive the conclusions
--- the correlation between standing ability and learning rate, and the
comparator's reliability --- are unmeasured working values, not estimates. They
are carried through every table that depends on them as explicit sensitivity
ranges rather than points, and the computations are structured so that revising
either is a one-line change in \path{scripts/compute.py} followed by a rebuild.
A reader who disputes either value can therefore recompute the report's numbers
under their own assumption rather than arguing about the conclusion.

% ==========================================================================
\section*{References}
\addcontentsline{toc}{section}{References}
\small
\sloppy

\begin{itemize}[leftmargin=1.2em,itemsep=0.25em,label={}]

\item AERA, APA, \& NCME (2014). \textit{Standards for Educational and
Psychological Testing}. American Educational Research Association.

\item Cohen, J. (1988). \textit{Statistical Power Analysis for the Behavioral
Sciences} (2nd ed.). Lawrence Erlbaum.

\item Dixon, C., Oxley, E., Gellert, A. S., \& Nash, H. (2023). Muddying the
waters: A systematic review of dynamic assessment of reading.
\textit{Reading and Writing}, 36(3), 673--698.
\href{https://doi.org/10.1007/s11145-022-10312-3}{doi:10.1007/s11145-022-10312-3}

\item Dorans, N. J. (2004). Equating, concordance, and expectation.
\textit{Applied Psychological Measurement}, 28(4), 227--246.

\item Faul, F., Erdfelder, E., Buchner, A., \& Lang, A.-G. (2009). Statistical
power analyses using G*Power 3.1. \textit{Behavior Research Methods}, 41,
1149--1160.

\item Gelman, A., \& Stern, H. (2006). The difference between
``significant'' and ``not significant'' is not itself statistically
significant. \textit{The American Statistician}, 60(4), 328--331.
\href{https://doi.org/10.1198/000313006X152649}{doi:10.1198/000313006X152649}

\item Haberman, S. J. (2008). When can subscores have value?
\textit{Journal of Educational and Behavioral Statistics}, 33(2), 204--229.
\href{https://doi.org/10.3102/1076998607302636}{doi:10.3102/1076998607302636}

\item Haberman, S. J. (2008). Subscores and validity.
\textit{ETS Research Report Series}, 2008(2).
\href{https://doi.org/10.1002/j.2333-8504.2008.tb02150.x}{doi:10.1002/j.2333-8504.2008.tb02150.x}

\item Hanley, J. A., \& McNeil, B. J. (1982). The meaning and use of the
area under a receiver operating characteristic (ROC) curve.
\textit{Radiology}, 143(1), 29--36.
\href{https://doi.org/10.1148/radiology.143.1.7063747}{doi:10.1148/radiology.143.1.7063747}

\item Hanley, J. A., \& McNeil, B. J. (1983). A method of comparing the areas
under receiver operating characteristic curves derived from the same cases.
\textit{Radiology}, 148(3), 839--843.
\href{https://doi.org/10.1148/radiology.148.3.6878708}{doi:10.1148/radiology.148.3.6878708}

\item Holland, P. W., \& Dorans, N. J. (2006). Linking and equating. In R. L.
Brennan (Ed.), \textit{Educational Measurement} (4th ed., pp. 187--220).

\item Kane, M. T. (2013). Validating the interpretations and uses of test
scores. \textit{Journal of Educational Measurement}, 50(1), 1--73.

\item Killip, S., Mahfoud, Z., \& Pearce, K. (2004). What is an
intracluster correlation coefficient? Crucial concepts for primary care
researchers. \textit{Annals of Family Medicine}, 2(3), 204--208.
\href{https://doi.org/10.1370/afm.141}{doi:10.1370/afm.141}

\item Livingston, S. A., \& Lewis, C. (1995). Estimating the consistency and
accuracy of classifications based on test scores.
\textit{Journal of Educational Measurement}, 32(2), 179--197.

\item MacCallum, R. C., Zhang, S., Preacher, K. J., \& Rucker, D. D. (2002).
On the practice of dichotomization of quantitative variables.
\textit{Psychological Methods}, 7(1), 19--40.
\href{https://doi.org/10.1037/1082-989X.7.1.19}{doi:10.1037/1082-989X.7.1.19}

\item Messick, S. (1995). Validity of psychological assessment.
\textit{American Psychologist}, 50(9), 741--749.
\href{https://doi.org/10.1037/0003-066X.50.9.741}{doi:10.1037/0003-066X.50.9.741}

\item Lohman, D. F. (2005). An aptitude perspective on talent identification.
\textit{Gifted Child Quarterly}, 49(2), 111--138.
\href{https://doi.org/10.1177/001698620504900203}{doi:10.1177/001698620504900203}

\item Meng, X.-L., Rosenthal, R., \& Rubin, D. B. (1992). Comparing correlated
correlation coefficients. \textit{Psychological Bulletin}, 111(1), 172--175.

\item NWEA (2025). \textit{MAP Growth Technical Report}. NWEA.

\item Ozen, Z., et al. (2024). A meta-analysis of the criterion-related validity
of the Cognitive Abilities Test. \textit{Gifted Child Quarterly}.
\href{https://doi.org/10.1177/00169862241285593}{doi:10.1177/00169862241285593}

\item Peters, S. J., Rambo-Hernandez, K., Makel, M. C., Matthews, M. S., \&
Plucker, J. A. (2019). Effect of local norms on racial and ethnic representation
in gifted education. \textit{AERA Open}, 5(2).

\item Sackett, P. R., Zhang, C., Berry, C. M., \& Lievens, F. (2022).
Revisiting meta-analytic estimates of validity in personnel selection:
Addressing systematic overcorrection for restriction of range.
\textit{Journal of Applied Psychology}, 107(11), 2040--2068.
\href{https://doi.org/10.1037/apl0000994}{doi:10.1037/apl0000994}

\item Sinharay, S. (2010). When do subscores have added value? Criteria based on
classical test theory. \textit{Journal of Educational Measurement}, 47(2),
150--174.

\item Taylor, H. C., \& Russell, J. T. (1939). The relationship of validity
coefficients to the practical effectiveness of tests in selection.
\textit{Journal of Applied Psychology}, 23(5), 565--578.

\item Thorndike, R. L. (1949). \textit{Personnel Selection: Test and Measurement
Techniques}. Wiley.

\item VanderWeele, T. J., \& Ding, P. (2017). Sensitivity analysis in
observational research: Introducing the E-value.
\textit{Annals of Internal Medicine}, 167(4), 268--274.

\item Westfall, J., \& Yarkoni, T. (2016). Statistically controlling for
confounding constructs is harder than you think.
\textit{PLoS ONE}, 11(3), e0152719.
\href{https://doi.org/10.1371/journal.pone.0152719}{doi:10.1371/journal.pone.0152719}

\end{itemize}

\normalsize

\subsection*{Repository sources}

Claims about the instrument's current state are grounded in the project's own
records rather than external literature:
\path{docs/product/STAGE2_BANK_RECOVERY_MEASUREMENT.md} (recovery ladder,
contamination floor, indeterminate rates);
\path{docs/product/STAGE2_QUESTION_DESIGN.md} (block design and gates);
\path{docs/research/ASSUMPTIONS_AND_EVIDENCE.md} entries E-095, E-100,
E-101, E-200 (recovery, validation-data commitment, cohort sizes, estimator
defect);
\path{docs/governance/DECISION_LOG.md} entries D-030 and D-200 (learning-rate
reportability);
\path{docs/research/IN_HOUSE_COGNITIVE_TEST_RESEARCH.md} and
\path{docs/research/EVIDENCE_DOSSIER.md} (CogAT and MAP psychometric
properties);
\path{docs/interviews/CRYSTAL_MARTEL_INTERVIEW_NOTES.md} (cohort sizes,
seat rationing, validation-data availability).


\end{document}
